\part[A diáspora judaica e o lugar]{A diáspora judaica\\ e o lugar}

\chapter{Fagulhas}

\noindent{}Em meio aos ofícios da Inquisição da Igreja Católica em 1492, em
Portugal e Espanha, contam as testemunhas que uma pessoa judia, se
negando a converter-se à força ao catolicismo, é levada à fogueira.
Junto a ela são colocados livros, folhas sagradas, pergaminhos e rolos
da Torá para que sejam consumidos pelas chamas, junto à sua
carne. A história nos conta que, no momento em que as primeiras
labaredas incendiaram os papéis, as finas letras timbradas do alfabeto
hebraico desprenderam-se, subindo em direção aos céus\footnote{Segundo
  as histórias de Rabi Uri de Strelisk: ``As miríades de letras da
  Torá correspondem às miríades de almas judaicas. Se faltar uma letra
  no rolo da Torá, a Torá não será válida; se faltar uma alma na união
  judaica, a Shechiná¹ não há de pairar sobre ela. Como as letras,
  também as almas devem unir-se e formar uma só união. Por que, porém,
  se proíbe que uma letra da Torá toque a outra? Porque toda alma
  judaica deve ter sua hora em que possa ficar a sós em conexão
  divina''. \emph{Shechiná} é uma palavra hebraica que se refere, de forma
  resumida, à figura de Deus no feminino e sua habitação e presença no
  plano terrestre. \textsc{buber}, \emph{História do Rabi}, p.\,464.}. Deste registro
flamejante de uma destruição que dispersa tudo --- coisas, letras,
corpos, histórias, almas, espaços, lugares, tempos ---, desses múltiplos
estilhaços --- de livros e de carnes --- outra coisa pode (re)nascer.
Que ``deste reduzir às cinzas nasça o seu próprio dar à
luz''\footnote{\textsc{didi-huberman}, \emph{Esparsas: Viagens aos papéis do
  gueto de Varsóvia}, p.\,126.}.

\begin{quote}
E vos digo: é preciso ter ainda caos dentro de si, para poder dar à luz
uma estrela dançarina. E vos digo: ainda há caos dentro de
vós\footnote{\textsc{bey}, \emph{\textsc{taz} --- Zona Autônoma Temporária}, p.\,7.}.
\end{quote}

Em março de 1943, no gueto de Vilnius\footnote{O gueto de Vilnius,
  na Lituânia, durou entre os anos de 1941
  e 1943.}, o poeta Avrom Sutzkever escreve: 

\begin{verse}
Talvez essas palavras também\\
Vão durar, e chegada a hora,\\
Subir para a luz\\
E florescer inesperadamente\\
E tal como o antigo grão\\
Tornou-se espiga\\
Talvez estas palavras vão nutrir,\\
Talvez essas palavras vão permanecer\\
Ao povo, em seu incessante caminho.\footnote{\textsc{sutzkever} \emph{apud}
  \textsc{didi-huberman}, \emph{op.\,cit.}, p.\,128.}
\end{verse} 

  No significado dessas palavras e
dessas imagens de sobrevivência, tais papéis-testemunhos, mais do que a
violência, perduram no tempo. A materialidade frágil de uma folha pode
até ser amassada, rasgada, transformada em cinzas, mas são nelas, nas
palavras --- escritas ou faladas ---, que as histórias permanecem vivas,
que resgatam sentidos profundos de resistência, atualizadas em cada uma
de suas gerações. E assim, fazem destes momentos de destruição um
``{[}\ldots{}{]} dar à luz a uma estrela dançarina''\footnote{\textsc{bey},
  \emph{\textsc{taz} --- Zona Autônoma Temporária}, p.\,7.}, um esforço para o
surgimento de desejos latentes de semear --- uma natureza de migalhas a
serem coletadas, e por fim espalhadas, como finos lapsos de fagulhas e
centelhas.

Desde as cinzas das chamas das fogueiras inquisitórias aos papéis dos
guetos, são estes os registros das odisseias particulares e coletivas,
de corpos individuais e comunitários que percorrem, transitam e ecoam
nos movimentos elipsoidais das diásporas, em que nascem e se constroem
territorialidades temporárias passíveis de apropriação e de
pertencimento. E assim, suplica-se: é hora de observar, de perto e de
longe, o que dessas outras localidades brotará\footnote{Pensar um
  judaísmo por uma ótica que considere as particularidades de uma
  identidade, de uma cultura e de uma localidade que se encontra (neste
  momento) em um novo processo de adaptação, resiliência, câmbio e
  debate interno.}. \emph{Mir zaynen do}\footnote{Do iídiche,
  ``Nós estamos aqui''. Expressão que dá voz ao hino dos
  \emph{Partisans}.}.

\begin{center}\adforn{49}\end{center}

\noindent{}Disponho três folhas sulfite, uma ao lado da outra, sobre a mesa de
jantar da sala de minha casa, rentes à beirada da tábua de madeira, em
uma falha tentativa de criar um fundo branco. Sobre as folhas estão dois
livros e um caderno. Registro esse primeiro passo meticulosamente
realizado. Em seguida posiciono a tão esperada fotografia preta e
branca, com alguns tons amarronzados, sobre os livros, o caderno, as
folhas e a mesa. Procuro por algo que chame minha consciência desde uma
certa distância --- temporal e espacial ---, uma ética de memória
situada para além do binômio vida-morte\footnote{O binômio vida-morte se
  estabelece no momento em que se lê as memórias trágicas do Holocausto
  (e demais tragédias --- neste caso em específico, tragédias judaicas)
  unicamente pela ótica da ``sobrevivência'' ou da ``não
  sobrevivência''. Neste tipo de leitura, nós --- judeus pós-Segunda
  Guerra Mundial (aqueles que sobreviveram --- nossa geração, nossos
  pais e nossos filhos) --- nos colocamos, sem perceber ou até mesmo sem
  querer, dentro de um espectro óptico traumático em que as memórias do
  Holocausto cotidianamente nos atormentam e assombram. Afinal, não são
  aqueles que sobreviveram que revisitam psiquicamente as situações dos
  guetos, dos campos de concentração e de extermínio? Não somos nós,
  sobreviventes, aqueles que pesquisam, debatem, choram e buscam ---
  mesmo sem fim aparente --- respostas para explicar tais tragédias?
  Neste texto procura-se entender que relembrar nossas dores e expor
  nossas fragilidades é defender a memória, e assim, a vida.}.

Estou diante do retrato. Nele uma jovem mulher posa com sua cabeça
levemente inclinada, exibindo por completo a lateral esquerda de seu
afinado rosto. De bochechas rosadas e olhos esverdeados, seus cabelos de
tons castanhos escurecidos criam um contraste com sua pele clara
destacada contra o acinzentado do plano de fundo da imagem.

A mulher retratada é a minha trisavó Dvora Najzsteter\footnote{Mãe
  biológica de meu bisavô Aba (Alberto) Najzsteter \textsc{z''l}, que por
  sua vez é pai de minha avó materna Sarah Dora Najzsteter Haber
  \textsc{z''l} (1941--2021).}, nascida em algum vilarejo nas proximidades
da cidade de Lublin, na Polônia, e que agora, já em minhas mãos,
lança-se em uma superfície fotográfica fosca que, apesar do ligeiro
sorriso, exprime a crueldade de uma história de que pouco ou nada
saberei. Surge uma série de questionamentos. Não seria a fotografia a
aposta de um testemunho? Não seria a fotografia a possibilidade de
contato e de distância ao mesmo tempo? Cabe a nós não ter medo de olhar.

Neste lapso de tempo congelado, que data entre as décadas de 1930 e
1940, Dvora veste um blazer azul-petróleo com seis finas linhas brancas
nas golas, onde se destaca um broche prateado. Seus cabelos escuros,
curtos e ligeiramente ondulados, são fixados com esmero por finos
grampos de tonalidade preta, que poderiam passar despercebidos a um
olhar menos cuidadoso. Seus olhos turvos se desviam da lente fotográfica
em um gesto aparentemente proposital de desencontro. 

E finalmente, será que de todas as coisas que poderiam estar olhando para ela agora,
ela imaginou por algum instante que seria eu?

No canto direito da fotografia, pelo olhar do espectador, está um
pequeno selo circular e timbrado que ocupa, aproximadamente, um quinto da
folha impressa. No centro se lê a palavra México, acompanhada dos
escritos: \textsc{foto herman e sta. maria de la redonda}, 136. Esse endereço não
existe mais. No verso do retrato, abrindo alas em meio aos vestígios do
tempo, encontra-se um pequeno texto, desejo de testemunho lido em
iídiche\footnote{O iídiche é uma língua judaico-germânica
  falada por judeus provenientes da Europa Central e Oriental (judeus
  ashkenazitas), uma mescla entre o alto-alemão do século \textsc{xiv} e
  elementos do hebraico e das línguas eslavas. Atualmente, apesar do
  grande número de falantes do iídiche serem judeus de vertentes
  ortodoxas e ultraortodoxas, o iídiche não é, em si, uma língua
  ortodoxa, mas sim uma língua diaspórica herdada e atravessada de
  espiritualidade e de tradição.} e escrito em letras hebraicas legíveis
e bem centralizadas na página, uma comprovação e, ao mesmo tempo, uma
palavra de adeus\footnote{Dvora fugiu da Polônia em torno do ano de
  1925, deixando o seu único filho --- órfão de pai --- aos cuidados de
  seus antigos sogros Dawid Najzsteter e Ruchla Lea Najzsteter. Ambos,
  Dawid e Ruchla Lea \textsc{z''l}, morreram em 1939 durante o início dos
  bombardeios nazistas nas cercanias da cidade de Lublin, na Polônia.}:
``Essa é uma fotografia minha de lembrança para os meus queridos
filhos. A tua mãe, Dvora''.

Sobre estes escritos --- ainda no verso do retrato --- encontra-se um
outro texto. Possivelmente uma primeira tentativa de aproximação entre a
autora, o papel e a caneta, mas pouco legível. As frases estão riscadas
e as letras se mesclam entre os traços azulados da caneta esferográfica.
As palavras parecem não possuir nem começo e nem fim.

% Comentário da editora: A preparadora trocou "este risco", que também não entendi a utilidade, por "os garranchos". São de fato garranchos? Não me incomoda, mas soa um pouco pejorativo pra um parágrafo que fala tanto sobre algo tão íntimo e pessoal. Também não entendi o trecho que segue, acho que precisa ser melhor explicado. A relação (sentimento, gancho com a pesquisa) com a fotografia da avó não está clara.

Os garranchos contrastam com a tonalidade já amarelada do suporte
fotográfico fosco, que me faz refletir. Penso sobre o apagamento da
lembrança, feito pela mão da própria autora --- e que ainda nestas
condições são lidas, desobedecendo o destino de uma morte prematura. São
gestos que não terminam, porque persistem no tempo. Como diz o poema de
Sutzkever: ``Talvez essas palavras também/\,Vão durar {[}\ldots{}{]}/\,Talvez essas palavras vão permanecer''\footnote{\textsc{sutzkever}
  \emph{apud} \textsc{didi-huberman}, \emph{op.\,cit.}, p.\,128.}. Já a segunda questão, o risco
da escrita por si só --- do ato de escrever letras que comprovam a
existência de um corpo, de uma pessoa em uma página. Uma mão, atrelada a
um corpo, escreveu essas palavras, que juntas se tornam frases,
finalmente são lidas enquanto desejo. \emph{Seria este o ritmo do fim?}
Entre essas promessas ardentes, prontas para serem consumidas pelas
chamas, estão também os papéis da Inquisição.

Ao olhar novamente percebo que, de minha trisavó Dvora, me restam
vestígios de sua caligrafia --- e assim uma terceira questão sobre o
risco surge\footnote{Proposta de risco como consequência.}:

% Comentário da editora: Eu acho que o iídiche deveria vir na nota, já que sempre priorizamos a acessibilidade do público (em português).

\begin{quote}
\emph{Ot di tokhike sakone tsu bashteyn oyf shraybn verter oyf yidish, vi di
vos shteyen ayngekritst in matn papir, vos ikh hob do far di oygn. Di-o
verter aleyn zaynen der konkreter bavayz fun iberlebn di natsishe
khayoles, fun firn a yidish lebn in der umgezetslekhkayt --- azoy az zey
zoln blaybn oyf yener zayt, bay der grenets tsvishn lebn un gezets. Azoy
arum az, in yenem moment, mer vi a shverd, antplekt zikh dos yidish fun
mayn elter-elter-bobe vi a pantser fun vidershtand kegn a
lingvistik-rasisher zoyberung, vos iz adurkhgefirt gevorn in eyrope. Dos
vayl, vi di yidish-meksikanishe shrayberin Myriam Moscona undz dermont:
\emph{Azoy vi mentshn, zaynen di shprakhn oykh farshikt gevorn in di
gazkamern}}\footnote{Este risco inerente de continuar a escrever
  palavras em iídiche como as que estão gravadas na folha fosca
  que agora tenho diante de meus olhos. Essas palavras são, por si só, a
  prova concreta de sobrevivência às tropas nazistas, à vida judaica na
  ilegalidade --- fazendo com que ela se mantenha no além, no limiar da
  vida e da Lei. E assim, naquele momento, mais do que uma espada, o
  iídiche de minha trisavó se mostra enquanto um escudo de
  resistência a uma higienização linguístico-racial em curso na Europa.
  Isso porque, como nos lembra a escritora judia-mexicana Myriam
  Moscona, ``assim como as pessoas, as línguas também foram
  enviadas às câmaras de gás''. \textsc{moscona}, \emph{Palabra del autor}.}
\end{quote}

Escrever em iídiche, em um gesto desesperado que vai de encontro
à fotografia, afugenta e retorna a mim. Como uma
resposta-eco\footnote{Trabalhar a experiência do eco como uma
  ressonância¹, algo que vai de encontro e que retorna em um duplo
  movimento de apropriação e distância. Pensar a ressonância enquanto tempo-limiar de retorno das
  perguntas que são colocadas no mundo. Ou seja, entre a primeira e a
  segunda aparição do signo, ou melhor, entre a ida e a volta de uma
  pergunta e/\,ou palavra ecoada, corre o tempo.}. Dessas memórias
permanece a \emph{nahalá}\footnote{Ritual judaico de reza e lembrança
  anual de entes queridos falecidos.}. E destas histórias de destruições
entoamos o \emph{Kadish}\footnote{Reza judaica de lembrança e elevação
  de alma e de espírito de entes falecidos em que há a santificação e a
  glorificação do nome de Deus.}. ``A nossa história não será em
vão''\footnote{Registro do Arquivo Ringelblum, acervo pertencente ao
  grupo judaico clandestino \emph{Oyneg Shabes} (em português, ``a
  alegria do \emph{shabat}''), que atuava coletando e registrando a vida
  comum e comunitária dentro do gueto de Varsóvia, entre os anos de
  1939--1943.}, suplicam todos aqueles judeus e judias que sobreviveram
ao desastre da vida judaica europeia entre as décadas de 1930 e 1940, e
continuam: ``todos saberão o que foi feito com os judeus na metade do século \textsc{xx}''\footnote{Registro do Arquivo
  Ringelblum.}. O iídiche, enquanto escudo, regressa para colocar
em xeque --- e assim, multiplicar --- tudo aquilo que por muito tempo
virou cinzas\footnote{A sobrevivência de uma língua diaspórica (como o
  iídiche e/\,ou o ladino) surge enquanto resposta às
  chamas flamejantes das fogueiras inquisitoriais e dos fornos
  crematórios dos campos de extermínio. São as próprias fagulhas das
  cinzas, esparsas partículas de papéis e de corpos que, lançadas para o
  ar e espalhadas pelo vento, se estabelecem e consolidam
  territorialidades próprias, dando sentidos de continuidade. Tanto o iídiche quanto o ladino são línguas
  judaicas diaspóricas. A primeira se constrói por meio de mesclas entre
  o alemão, o hebraico e alguns elementos de línguas eslavas. E a
  segunda por meio de mesclas entre o português e o espanhol medieval, o
  italiano e o hebraico. O iídiche é o idioma comum falado por
  judeus ashkenazitas (judeus provenientes da Europa Central e do
  Leste Europeu) enquanto o ladino é o idioma falado por judeus
  sefaraditas (judeus provenientes de Portugal e Espanha).}.

  % Comentário da editora: É só "metade" mesmo? Tem "primeira" ou algo assim? Seria necessário colocar [primeira] assim, entre colchetes, para indicar algo? Também acho que até dá pra entender o que "coloca em xeque" (que seria a própria validade da língua e da cultura iídiche), mas talvez faça mais sentido excluir essa frase. Isso foi algo que a preparadora apontou. Ou então reescrever essa parte.

Toda sobrevivência --- seja ela construída pelos seus maiores ou pelos
seus menores gestos --- é uma espécie de contestação histórica. Toda
sobrevivência é um desejo pulsante e latente de escrever\footnote{Para a
  maioria das pessoas judias perseguidas (mortas em bombardeios, dentro
  dos guetos, em campos de extermínio ou em valas comuns) não houve
  milagre. Suas vidas não foram salvas, para alguns apenas as suas
  palavras foram e, assim, persistem no tempo fazendo com que suas vidas
  e histórias sejam lidas ainda hoje, depois de suas mortes.}, de
registrar essa breve passagem, de fazer parar a roda da História. Assim
como os levantes, como diria o escritor e filósofo Hakim Bey\footnote{Em
  seu livro \emph{\textsc{taz} --- Zonas Temporárias Autônomas}, o escritor e
  filósofo anarquista Hakim Bey descreve a propagação de espaços
  autônomos temporários enquanto táticas de resistência e de
  esvaziamento de poder. Não seria o próprio ato de resistência do grupo
  clandestino judaico \emph{Oyneg Shabes} um espaço autônomo? Não seriam
  as páginas escritas e guardadas, para posteriormente serem resgatadas,
  criadoras de espacialidades temporárias? São possibilidades de outros
  mundos, e assim persistem no tempo.}, as tentativas de sobrevida são
insurreições que surgem para além do Tempo e violam a lei da História.
Os levantes --- em seus maiores ou em seus menores gestos --- e a
libertação de uma territorialidade --- em qualquer escala ---
possibilitam movimentos para fora e para além da \emph{espiral
hegeliana} de progresso. É dessa maneira que entendo as ligeiras
palavras em iídiche de minha trisavó --- um ato de registrar sua
breve passagem e gravar seus anseios de adeus e de amor.

Assim como nos contos talmúdicos, em que Deus criou o mundo (e
todas as coisas que ele contém) por meio da codificação das 22 letras
dispostas do alfabeto hebraico, o iídiche de minha trisavó se
torna --- naquele momento --- seu mundo por completo. E finalmente,
neste imenso oceano de apagamentos e refazeres, talvez --- agora ---
seja a hora de partir novamente do começo, e dessa vez com mais
distância.

\begin{center}\adforn{49}\end{center}

\noindent{}Existe uma história talmúdica que diz que o mundo foi criado e
destruído 26 vezes. O mundo já nasceu da e na ruína. Rabi
Bunam\footnote{Rabi Bunam de Pzhysha (1765--1827) foi um dos principais
  líderes do movimento \emph{hassídico} na Polônia durante o final do
  século \textsc{xviii} e o início do século \textsc{xix}. Seus contos, provérbios e
  histórias foram coletados e compilados no livro \emph{História do
  Rabi} do historiador e filósofo Martin Buber.} costumava dizer:
``No dia do Ano Novo\footnote{\emph{Rosh Hashaná} (em português,
  ``Cabeça do Ano'') é o dia festivo de início do calendário judaico.
  Acontece no primeiro dia do mês judaico de Tishrei, que corresponde
  normalmente a finais de setembro e inícios de outubro no calendário
  gregoriano.} o mundo começa novamente e, antes de começar,
termina. Assim como, antes de morrer, todas as forças agarram-se à vida,
o homem, no dia do Ano Novo, deve agarrar-se à vida do mundo com todas
as forças''\footnote{\textsc{buber}, \emph{op.\,cit.}, p.\,577.}.

A ``ruína'' e o ``agarrar-se à vida'' que nestas páginas
tenta-se construir --- por meio do esmiuçamento das territorialidades
temporárias, das línguas, da diáspora --- estão presentes nos documentos
e nos arquivos coletados e armazenados pelo grupo judaico \emph{Oyneg
Shabes}\footnote{\emph{Oyneg Shabes} foi um grupo judaico clandestino
  liderado pelo historiador polonês Emanuel Ringelblum (1900--1944) que
  atuou entre os anos de 1939 e 1943 dentro dos limites do gueto de
  Varsóvia, Polônia. O grupo dedicava-se a coletar e preservar relatos,
  histórias, sonhos, contos, gritos, choros e demais testemunhos sobre a
  vida individual e comunitária no gueto de Varsóvia durante a Segunda
  Guerra Mundial. Algumas horas antes do levante do gueto, em 1943, o
  grupo enterrou seus arquivos em um ato de esperança de que esses
  documentos fossem encontrados ao final da guerra e que suas memórias
  pudessem, ao menos, ser preservadas.} --- que dedicava-se a escavar
cada centímetro da história vivida entre 1939 e 1943 dentro do gueto de
Varsóvia. Seus registros são uma recolha incansável e imensurável de
vida (e de morte) e têm por intuito conservar migalhas de histórias,
mesmo daquelas mais particulares\footnote{A noção de particulares, que
  aqui tenta-se construir, está ligada ao entendimento de que as
  famílias, as comunidades e os indivíduos dentro do gueto possuíam suas
  próprias histórias e trajetórias de vida. Mesmo assim, suas histórias
  singulares, ao final, são formadas por meio de papéis amontoados que
  ressoam, por mais diferentes que sejam, em uma história coletiva.}.

Em \emph{Cosmopoéticas do refúgio}, o filósofo Dénètem Touam Bona, sobre
a identidade fronteiriça e a tentativa de construção de mundo que, ainda
hoje, torce a ``linha da cor'', diz:

\begin{quote}
Sobreviver num universo em ruína é se deixar atravessar, deixar-se
habitar pelo acréscimo da vida prodigado pelos ancestrais, pelos
animados (animais, plantas e povos do infinitamente pequeno) e pelos
elementos: abraçar a própria morte, a potência da sombra e do húmus,
para renascer\footnote{\textsc{bona}, \emph{Cosmopoéticas do refúgio}, p.\,85-86.}.
\end{quote}

Os testemunhos do grupo \emph{Oyneg Shabes}, essa espécie de
``{[}\ldots{}{]} abraçar a própria morte {[}\ldots{}{]} para
renascer''\footnote{\emph{Ibidem}.}, que por longo tempo permaneceram sob os
escombros do gueto de Varsóvia --- soterrados de neve e de terra ---,
manifestam-se não somente enquanto uma busca por uma história
particular, mas também enquanto uma narrativa montada com papéis
desatados por cada uma das histórias e das tragédias singulares. Este é
o caso de minha trisavó Dvora e de milhares de outras pessoas fugitivas
de guerra, que deixaram suas famílias em busca de promessas de
vidas\footnote{E pensar, o que elas diriam sobre Gaza?}. Afinal, como
nos lembra Didi-Huberman\footnote{\textsc{didi-huberman}, \emph{op.\,cit.}, p.\,107.}, o
paradoxo está em entender que judeus e judias foram unidos em morte, mas
que, em vida, viveram esparsamente: atravessados por brechas e fissuras.

Na falta de outras verdades anseio em acreditar que os parentes de minha
trisavó estavam unidos, pelo papel e pela carne, entre aqueles judeus e
judias que se levantaram contra o regime nazista durante o levante do
gueto de Varsóvia em 1943; e quiçá suas memórias, seus
``papéis-desejos'', seus ``papéis quaisquer'' e seus
``papéis chorados'' possam, ao menos, ter sido coletados,
armazenados e enterrados pelo grupo \emph{Oyneg Shabes}. E que assim
suas memórias, seus registros, suas histórias --- mesmo que eu não os
possua --- possam, de alguma maneira, perdurar no Tempo.

Os cunhados-irmãos\footnote{Após a morte de seu marido, Dvora foge da
  Polônia em direção aos Estados Unidos da América, deixando seu filho,
  meu bisavô Aba (Alberto) Najszteter, aos cuidados de seus cunhados e
  de seus sogros. A expressão ``cunhados-irmãos'' quer traduzir a
  proximidade que Dvora teria com os irmãos e irmãs de seu falecido
  marido.} de minha trisavó Dvora\footnote{Aqui é a nossa obrigação
  dar-lhes nome: Moshe Mekhl Najszteter \textsc{z''l}, Beila Myriam
  Sherberg \textsc{z''l}, Icchak Leib Najszteter \textsc{z''l}, Yehudit
  Sheindla Sukholski \textsc{z''l}, Mordechai Yossef Najszteter
  \textsc{z''l}, Hanna Sifra Zilbershpan \textsc{z''l}, Tova Shifra Godnizki
  \textsc{z''l}.} foram sistematicamente enviados ao campo de concentração
e de extermínio de Treblinka, na Polônia, a uma hora e meia de trem da
cidade de Varsóvia. Assim como foram os mais de 300 mil judeus e judias
residentes no gueto de Varsóvia --- dos quais seus cunhados-irmãos
também faziam parte --- e que se levantaram contra o regime nazista
entre abril e maio de 1943. Dos testemunhos-registros dos
cunhados-irmãos de minha trisavó Dvora me restam alguns nomes,
sobrenomes e longínquos graus de parentesco --- possivelmente alguma
carga genética --- o fio genealógico\footnote{A preservação genealógica
  --- seja ela sanguínea, textual ou cultural --- é uma das mais
  importantes preservações da vida das comunidades judaicas nas
  diásporas. Segundo o escritor Amós Oz e a historiadora Fania Oz-Salzberger:
  ``Para ser judeu não é necessário ser \emph{filho/\,filha de mãe judia},
  basta apenas ser um leitor''. Isso porque, na tradição judaica, todo
  leitor é revisor de originais, todo aluno é crítico e todo escritor,
  inclusive o criador do Universo, incorre em um grande número de
  questões.}. Dos milhares de judeus residentes do gueto de Varsóvia nos
restam resquícios de histórias, sonhos, canções,
gritos, ``papéis-desejos'', promessas de retorno, de persistência,
de vida. Dos judeus e judias perseguidos e mortos durante o Holocausto
nós mantivemos --- conservamos e fazemos perdurar --- ínfimos impulsos
de sobrevivência, gestos de resistência, fugas, resiliências e
sobrevidas, hinos partisanos e ``a lamentação para que nos
ensine, nos levante''\footnote{\textsc{didi-huberman}, \emph{op.\,cit.}, p.\,129.}.

Hoje, Treblinka --- este campo despovoado\footnote{Os nazistas
  destruíram não apenas os lugares judaicos de culto, de reza, de vida e
  de comunidade, como também seus lugares de morte.} e de morte, sítio
arqueológico, e possivelmente um suposto memorial --- se preenche de
deformadas pedras por toda a sua extensão. São formações rochosas
gravadas --- em nome\footnote{A morte pode possuir muitos nomes. Nas
  mais diversas pedras, espalhadas ao longo do memorial-campo de
  Treblinka, estão gravados nomes de lugares de memória que jamais
  devemos nos esquecer.} e em carne --- numa terra infértil. São elas,
as próprias pedras, memórias eternas, faíscas das chamas das almas de
pessoas mortas, que, espraiadas pelo ar, e levadas pelo velho vento do
sopro nômade da diáspora, se tornam fagulhas em que suas recordações e
memórias persistem e perduram no tempo, permanecem e criam
territorialidades temporárias por meio da língua, da conservação da
tradição, do indeterminismo da diáspora, e assim, são levadas de geração
em geração. São as faíscas das almas as próprias memórias que retornam;
e como diria Rabi Bunam, que antes da ruína agarram-se à vida.

Na tradição judaica há o costume de assentar pedras em túmulos nos
cemitérios como manifestação de recordo, de presença e da crença na
memória permanente. Certamente a materialidade de uma pedra não é
infinita, porém quiçá seja ela eterna em um determinado espaço-tempo ---
e, assim como a complexidade de Deus, essa energia não é infinita porque
não tem fim, antes é eterna porque não tem começo\footnote{Como diria
  Gabriel García Márquez: ``É a vida, mais do que a morte, que não
  tem limites''.}.

Após a guerra, foram encontradas, ao todo, 35.369 páginas esparsas
enterradas nos subsolos do gueto de Varsóvia pelo \emph{Oyneg Shabes}.
Este total representa apenas dois terços do material enterrado pelo
grupo. O restante --- um terço do material --- jaz --- afogado em terra
e em lama --- nas profundezas fúnebres de um solo gélido. Curiosamente,
o ``escovar a História a contrapelo''\footnote{``Escovar a
  história a contrapelo'' é uma metáfora de Walter Benjamin que busca
  analisar o passado questionando e revisitando a narrativa oficial dos
  vencedores. Por meio dessa abordagem, foca-se nas experiências dos
  oprimidos, vencidos e silenciados, revelando as lacunas e resistências
  que o progresso linear ignorou.}, sobre que o filósofo Walter Benjamin
há tempos já discorria, é literalmente materializado nos escombros do
gueto. São pedras, tijolos, pedaços de concreto, estilhaços de vidro e
demais resíduos da construção civil levantados e retirados em busca de
testemunhos, de folhas sujas, amassadas e soterradas --- e que, por fim,
sobrevivem às intempéries da História e do Tempo.

% Comentário da editora: Acho que está tudo bem usar história e tempo em caixa baixa aqui. Existem também outras ocasiões em que termos parecidos estão em caixa alta sem necessidade. Posso padronizar?

Usando novamente da deposição do filósofo Dénètem Touam Bona: ``As
guerrilhas se apresentam como uma não batalha, um disfarce e uma
camuflagem que zombam com a moral dos poderosos''\footnote{\textsc{bona}, op.
  cit., p.\,41.}. Dessa maneira, o grupo \emph{Oyneg Shabes}, para além
de uma tentativa de preservação histórica, brilhantemente entende que,
para o ``povo do livro'', basta também que suas páginas
sobrevivam, que seus traços de humanidade possam ser lidos em letras
sobre papéis. É por meio do encontro entre o fino risco da caneta e a
matéria que a palavra transforma-se em lugar. É a partir destes gestos
que recordamos os 6 milhões de judeus e judias mortos durante o
Holocausto --- e assim, compreendemos que suas almas retornam e
retornarão, agarrando-se à vida --- perdurando em cada um dos finos
traços de letras.

Dedica-se este texto aos de identidade fronteiriça. E entende-se que
``toda identidade é incompleta, sem uma imagem de
alteridade''\footnote{\textsc{sorj}, \emph{Sociabilidade brasileira e identidade
  judaica,} p.\,21.}.

Que essas memórias sejam uma revolução. E ``que nenhum rosto da
realidade humana/\,seja empurrado para baixo do silêncio da
História''\footnote{\textsc{tansi} \emph{apud} \textsc{bona}, \emph{op.\,cit.}, p.\,3.}.

\begin{center}\adforn{49}\end{center}

\noindent{}Ainda prefiro a versão das coisas em que não se sabe ao certo o que foi
lido, o que foi inventado. Se diz que existem apenas sete possíveis
histórias a serem contadas no mundo, e o que muda entre cada uma delas
seria a capacidade de rearranjo de cada coletor, de cada recorte, de
cada ritmo, de cada jeito de entrelaçar essas narrativas. As sete
histórias não saberia dizer, porém sei que continuamos a escrever sempre
o mesmo livro.

Em \emph{A ensaificação de tudo}, a escritora norte-americana Christy
Wampole diz que ``O ensaísta está interessado em pensar sobre si
pensando sobre as coisas''\footnote{\textsc{wampole}, \emph{A ensaificação de
  tudo}.}. Ao fim e ao cabo, esta pesquisa é uma espécie de ensaio que,
assim como a própria diáspora, sustenta-se enquanto tentativa de
solução, ainda que na ausência de soluções. São impulsos e propostas
para ``pensar sobre si pensando sobre as coisas''\footnote{\emph{Ibidem}.},
e dessa maneira, construir outras territorialidades e frequências
comunitárias. Assim sendo:

\begin{quote}
Esse não é um livro de história. A escolha que nele se encontrará não
seguiu outra regra mais importante do que meu gosto, meu prazer, uma
emoção, o riso, a surpresa, um certo assombro ou qualquer outro
sentimento, do qual teria dificuldades, talvez, em justificar a
intensidade, agora que o primeiro momento da descoberta
passou\footnote{\textsc{foucault}, \emph{A vida dos homens infames}, p.\,203.}.
\end{quote}

Busca-se aqui entender que as mais diversas \emph{perspectivas
judaico-diaspóricas sobre o lugar} estão entremeadas e entrelaçadas de
signos, símbolos, línguas e arquiteturas --- em suas mais variadas
escalas --- em que se permite o nascimento e a consolidação de
identidades fluidas, regadas de indeterminabilidade e, por isso,
passíveis de apropriação e de pertencimento.

São as fagulhas, as chispas, as faíscas, as memórias-palavras que
perduram, que vagueiam pelo ar e são levadas pelo velho sopro do vento
nômade da diáspora. E, ao pousarem, criam territorialidades que, mesmo
que temporárias, possibilitam a construção de localidades próprias e de
sentidos de identificação. Dessa maneira, os lugares que aqui procura-se
trabalhar, mais do que espacialidades físicas, se manifestam em seus
sentidos identitários e comunitários e através das línguas, dos objetos
culturais e de ritos religiosos. São eles: os espaços de trânsito, do
indeterminado, do fluido e do híbrido que cruzam --- em suas mais
diversas maneiras --- as noções espaço-temporais judaicas.

Essas geografias de identidades apropriadas serão esmiuçadas e
investigadas no decorrer deste texto por meio dos conceitos judaicos de:

\begin{enumerate}
\item \textsc{bs'd}; 
\item Kipá; 
\item Mezuzá; 
\item Eruv; 
\item Makom. 
\end{enumerate}

Transversalmente, procura-se entender
a página, as palavras, as línguas diaspóricas\footnote{As línguas também
  possuem biografia. O ladino conserva as memórias de infância de
  um espanhol medieval. As pessoas podem esquecer, mas as línguas não.
  As línguas mantêm as alegrias, as dores e as feridas de tudo que há
  passado.} como elementos essenciais para a criação de corpos
identitários e comunitários; a \emph{Shechiná} --- o lugar de habitação,
de crença e de presença divina --- como uma territorialidade temporária;
e a diáspora enquanto um limite e um limiar em si, um espaço eternamente
transitório.

Rabi Tarfon\footnote{Rabi Tarfon foi membro da terceira geração de
  estudiosos da \emph{Mishná}, vivendo entre os séculos \textsc{i} e \textsc{ii} d.e.c.}
costumava dizer: ``Não lhe é exigido que complete a tarefa, mas
você não é livre para escapar dela''\footnote{\emph{Pirkei Avot} 2:16,
  escrito cerca do século \textsc{i}.}. Essa constante elaboração é o que
impossibilita o esgotamento de sentidos sobre o ser judeu. Por fim, já
que estaremos sempre a escrever o mesmo livro, essa investigação
sustenta-se enquanto um ensaio identitário, sem a pretensão de alicerçar
e fixar imagens judaicas.

Em resumo, esta breve, dispersa e pequena introdução é uma abertura de
portas.

\emph{I ke io lo mire}\footnote{Do ladino: ``e que eu o veja''.}.

\chapter{Páginas}

\noindent{}Ao final da tarde do sexto dia da criação divina --- \emph{erev shabat}
---, às vésperas do dia sagrado, de descanso religioso e espiritual,
alguns minutos antes do pôr do sol, Rabi Hanina e Rabi Oshea, dois
rabinos talmúdicos do século \textsc{iii}, se reuniam nas encostas do
Monte Moriá\footnote{O Monte Moriá é um território
  sagrado judaico e o local escolhido para a construção do primeiro e do
  segundo Templo Sagrado de Jerusalém¹. Possui um simbolismo cósmico e
  testemunhou uma das máximas práticas de benevolência judaica:
  \emph{Veahavtá Lereacha Kamocha} (``Ame o teu próximo como a ti
  mesmo'') --- acrescentado por Rabi Akiva: \emph{Ze'Klal Gadol baTorá}
  (``Este é o maior princípio da Torá''). A filosofia judaica acredita que o terceiro Templo Sagrado de
  Jerusalém --- bem como a salvação divina --- virá até o ano 6000 do
  calendário judaico (atualmente nos encontramos no ano 5786).}, em
Jerusalém. Nesse recinto sagrado, moldavam e esculpiam conjuntamente um
animal em barro e gravavam em sua testa a palavra
\emph{emet}\footnote{\emph{Emet}, além de significar em português
  ``Verdade'', é um dos nomes judaicos de Deus.}. Com o auxílio do
\emph{Sêfer Ietsirá}\footnote{O \emph{Sêfer Ietsirá} (em português: ``O
  Livro da Criação'') é um dos mais antigos tratados filosóficos
  escritos em língua hebraica.} traziam à vida um bezerro e, logo em
seguida, o cozinhavam como alimento para a ceia festiva.

Em \emph{Terra e Paz}, o escritor e poeta israelense Yehuda Amichai nos
lembra que a procura por um bezerro\footnote{No original, um cabrito. Ou
  qualquer outro animal de pasto.}, ou por um filho\footnote{O Monte
  Moriá foi o local em que o primeiro patriarca Abraão foi
  impelido a sacrificar seu filho Isaac a fim de provar a sua lealdade a
  Deus.}, sempre foi o começo de uma nova religião por estes
montes\footnote{\textsc{amichai}, \emph{Terra e Paz}, p.\,29.}.

\begin{center}\adforn{49}\end{center}

\noindent{}Conta o Talmud que, durante a época do segundo Templo Sagrado de
Jerusalém, uma jovem mulher, proibida de transitar por esta
territorialidade\footnote{Na maioria das comunidades judaicas, e em
  especial naquelas comunidades tradicionais, ortodoxas ou
  ultraortodoxas, mulheres ainda são proibidas de assumir papéis de
  liderança no âmbito religioso e comunitário.}, envolve-se em uma
conversa com Deus, por meio de palavras. A jovem suplica a Deus para que
conceda a ela um filho para criar, permitindo, desta maneira, que seja
portadora de uma nova linhagem de descendentes e matriarca destes fios
genealógicos. Assim, ao envolver-se em uma conversa com Deus, a jovem
mulher passa a ser, ela mesma, porta-voz de sua identidade, de sua
conexão e de sua relação pessoal com a divindade\footnote{No judaísmo
  pós-templo, qualquer pessoa pode se conectar com Deus e com a
  divindade sem a necessidade de uma intermediação e encontrar
  interpretações próprias nos textos sagrados.}.

Entre 586 a.E.C. e 76 d.E.C., período em que perdurou o segundo Templo
Sagrado de Jerusalém, os hebreus se comunicavam com a divindade por meio
de sacrifícios e oferendas realizadas e concedidas pelo \emph{Cohen
Hagadol} --- o Sumo Sacerdote de Israel. Este era o canal da comunicação
direta e unilateral com Deus.

Após a destruição do Templo Sagrado, a expulsão da
\emph{Shechiná}\footnote{\emph{Shechiná} é uma palavra hebraica que, de
  maneira resumida, significa ``habitação'' ou ``presença de Deus''.}, e
assim a inexistência e a insuficiência de um território próprio judaico
para a realização de oferendas e de sacrifícios, os judeus passam a
comunicar-se com a divindade por meio de livros, de textos e do
pergaminho. Deus e, em especial, a memória do Templo Sagrado\footnote{A
  memória do Templo Sagrado e a lembrança da cidade de Jerusalém
  permeiam as rezas, os ritos, as celebrações, as festividades, bem como
  os momentos de luto judaico. Todas as sinagogas devem ser construídas
  mirando o Templo Sagrado (uma Jerusalém espiritual), voltando as
  orações a esta localidade santa. Nos noivados judaicos há o costume
  das mães dos noivos quebrarem conjuntamente um prato, nos casamentos
  judaicos quebra-se um copo, nas casas judaicas há uma parede sem
  revestimento (apenas em tijolos). Todas essas ações remetem à
  incompletude e à ausência do Templo Sagrado.} passam, dessa forma, a
habitar o texto e o livro.

Para manter a continuidade judaica viva, este novo aspecto da
comunicação com a divindade é santificado por meio de rezas e de cultos
religiosos entoados coletivamente, nas primeiras sinagogas da diáspora.
As rezas e as orações judaicas são conversas com Deus. São poemas,
cânticos e bênçãos que reiteram a lealdade e a fidelidade ao monoteísmo
e à divindade judaica. A reza judaica é um sopro entoado e proferido em
voz alta (em hebraico) e em comunidades em que os indivíduos judeus se
unem --- formando um \emph{minian}\footnote{\emph{Minian} é o quórum
  mínimo de dez homens judeus adultos (pós-\emph{bar-mitzvá}, maioridade
  religiosa) necessários para realizar certas obrigações religiosas, bem
  como para a retirada da Torá da Arca Sagrada e sua leitura.
  Mais recentemente e em algumas vertentes religiosas e liberais
  judaicas, as mulheres também passaram a contar como participantes
  ativas do \emph{minian}.} --- para realizar os ritos, para rezar e
santificar o nome de Deus de maneira coletiva.

Segundo a narrativa bíblica da criação Deus moldou, à sua imagem, o
homem em barro (em hebraico, \emph{Adamá}), o qual insuflou com o Seu
sopro. \emph{Adam} (Adão), o primeiro ser humano, nasce do cruzamento da
terra, da mão de Deus que confeccionou e moldou o barro e do sopro
divino --- adicionando vida, espírito e uma alma a esta matéria amorfa.

Conta o Talmud que Deus criou o mundo a partir do ar primal e
gravou, nesta esfera invisível, as 22 letras disponíveis do alfabeto
hebraico. É por meio deste gesto que todas as coisas do mundo, todos os
seres dos três estratos do cosmo --- espaço, tempo e o corpo humano ---
passam a existir. Assim, ``as palavras existem porque correspondem
a coisas existentes, e as coisas deveriam existir porque há palavras
para nomeá-las''\footnote{\textsc{manguel}, \emph{A musa da impossibilidade}, p.\,42.}. Por fim, o encadeamento linear das 22 letras hebraicas se torna
palavras que nomeiam coisas e, finalmente, compõem frases e orações.

No momento da Inquisição, entre as entranhas da espadana do resquício do
fogo, em que as inscrições das palavras sagradas se desprendem, voando
em direção aos céus, uma imagem assemelha-se ao ar primal de Deus. Tais
letras, ao se desgrudarem dos papéis, superam a própria morte, pousam em
outros territórios, criam mundos e fixam-se nestes espaços de vácuo. Ou
seja, as letras, ao morreram para renascer, tornam estes espaços lugares
de apropriações comunitárias e reiteram a criação do mundo por Deus.

Assim como o verbo ``escrever'' (em latim, \emph{scribere}), um gesto de
fazer uma incisão sobre um objeto e/\,ou corpo, o homem foi inscrito e
criado a partir do barro, ou seja, gravado e inserido nele. Ao final,
``os seres humanos'', como menciona Flusser, ``são
materiais espiritualizados da fusão da terra com Deus''\footnote{\textsc{flusser},
  \emph{A escrita. Há futuro para a escrita?}, p.\,25.}. Ao inscrever o
homem no barro, Deus grava, fura e penetra uma materialidade --- a molda
---, se apropria de sua maleabilidade física para criar o homem e,
assim, trazê-lo à vida.

Deus insere nesta matéria amorfa e desprovida de inteligência própria
--- um \emph{golem} --- uma alma por meio do sopro. Talvez essa seja
também a ação realizada pela jovem mulher. Ao rezar, ao iniciar um
diálogo com Deus, a jovem sopra palavras que são proferidas, voam e
vagueiam pelo ar até ganharem vida --- uma alma ---, até serem
espiritualizadas. Essas palavras, na forma de orações, flutuam em
direção aos céus, rompem os portais celestiais divinos, permitindo que a
divindade seja finalmente acessada sem uma intermediação.

Como diz a escritora e poeta brasileira Marília Garcia, ``as falas
acontecem no presente''. As falas e orações que saem de nossas bocas
por meio de vultos e frestas entre a língua e o ar da linguagem são
todas entoadas e ressoadas no momento do agora.

As falas, as palavras, as comunicações, bem como as histórias
testemunhadas e transmitidas oralmente (ou de forma escrita) de geração
em geração, acontecem no momento presente e, assim, se adiciona,
acrescenta e multiplica a alma, dando vida às comunidades, aos ritos e
aos cultos judaicos. Estes, por fim, também são interpretados de modos
diferentes, segundo o momento histórico, adicionando camadas de
complexidade às tradições.

Ao soprar, Deus inseriu uma alma e um espírito nos primeiros seres
humanos --- Adão e Eva --- e os contemplou com a vida. Ao final,
``o espírito penetra o objeto por ocasião do inscreve, para \emph{dar-lhe-vida}''\footnote{\emph{Ibidem}, p.\,25.}. Dessa forma, ao inscrever
uma informação (uma alma) em uma matéria amorfa, o autor --- Deus na
criação do mundo e de todas as coisas que ele contém, ou os seres
humanos por meio de frases e de palavras\footnote{Posso não acreditar em
  Deus, mas acredito no poder das palavras --- e, para mim, é quase a
  mesma coisa.} --- deseja que estes objetos e/\,ou corpos, por ele
inscritos, ressoem, perdurem no espaço e no tempo e, ao menos, não se
decomponham tão rapidamente.

% Comentário da editora: Tem algo errado na citação acima.

Estes objetos e/\,ou corpos inscritos e escritos --- possivelmente lapsos
de fagulhas ---, ou melhor, a vontade de perdurar e perfurar o espaço e
o tempo, se fazem sentir nos papéis do grupo judaico \emph{Oyneg Shabes}
como um nítido desejo de vida e de persistência. E assim também nos
escritos de minha trisavó Dvora. Uma vontade de adicionar --- a todo
custo --- a estas folhas, a estas páginas e a estes documentos (e a
estas palavras) uma vida própria, uma alma e um espírito que fure o
tempo, para que não nos esqueçamos tão rapidamente.

Em \emph{Teste de Resistores}\footnote{GARCIA, \emph{Teste de
  Resistores}, p.\,14.}, Marília Garcia comenta o diálogo entre Ferdinand
e Marianne em \emph{Pierrot le Fou}, filme de Jean-Luc Godard. Após um
assassinato e o roubo de um carro, Ferdinand quer parar para descansar
em uma praia. Ele se diz apaixonado, mas Marianne parece pensar apenas
em dinheiro.

\begin{verse}
nesse momento ferdinand se vira\\
olhando para trás na direção da câmera\\
e diz --- estão vendo\\
ela só pensa em se divertir\\
marianne se vira também olhando para trás\\
na direção da câmera\\
e pergunta para ele --- com quem você está falando?\\
ao que ele responde --- com o espectador.
\end{verse}

Este pequeno diálogo contribui para dar ao filme a dimensão de filme ---
apresentando o furo do cinema. Há pessoas que acreditam que não poderiam
viver sem escrever, \emph{Oyneg Shabes} e minha trisavó Dvora escrevem e
``quem escreve ultrapassa o ponto final em direção ao outro, o
leitor''\footnote{\textsc{costa} \emph{apud} \textsc{flusser}, \emph{op.\,cit.}, p.\,8.}. Essas pessoas
acreditam que o gesto de escrever pode expressar e atestar as suas
existências\footnote{Eu também escrevo porque as palavras me pertencem,
  e assim reafirmo a minha fé no poder da linguagem.}. Assim,
\emph{Oyneg Shabes} e minha trisavó estendem a mão em direção ao outro a
fim de provocar uma inquietação; comunicam-se com seus interlocutores e
permitem que suas palavras escritas falem por si.

Por fim, supõe-se que a inscrição (a criação) feita por Deus, dos
primeiros seres humanos, é repetida simbolicamente todos os anos, já que
todos os anos o mundo é novamente cosmogonizado por Deus e a aliança é
reafirmada. Segundo a filosofia judaica, Deus, no momento de \emph{Iom
Kipur}\footnote{O \emph{Iom Kipur} é uma das mais importantes
  festividades do calendário judaico, momento em que a proximidade de
  Deus permite a purificação dos pecados diante da divindade e com o
  próximo.}, inscreve e sela (outra vez) a humanidade em mais um ano de
vida, reitera a confiança na criação divina do mundo, e assim expande
novamente o sopro, a alma e uma vida.

Essa afirmação da aliança, e da lealdade judaica com Deus por meio da
palavra, é também realizada durante o momento do \emph{brit-milá} (em
português, ``aliança da circuncisão''). O \emph{brit-milá} é uma
cerimônia ritualística em que o prepúcio de meninos no oitavo dia de
vida é cortado, como símbolo da lealdade entre Deus e o povo judeu. É
neste momento, também, que o menino judeu recebe o seu nome diante da
comunidade --- e, assim, passa a existir perante ela e a frequentá-la.

Em hebraico, \emph{milá} --- a depender de sua forma escrita --- pode
significar tanto ``circuncisão'' ({מילה}), quanto ``palavra''
({מלה}); dessa maneira, uma tradução possível deste ritual (que
envolve um recém-nascido --- uma nova geração de indivíduos judeus)
seria: ``aliança da palavra''. Assim, a lealdade judaica a Deus e a
crença nas palavras são reafirmadas nos corpos judaicos por meio de uma
aliança com a lei judaica escrita, com a Torá, com os textos, com
as tradições e com a divindade --- e com a comunidade.

Se as falas, as rezas e as bênçãos são entoadas enquanto sopros, e
estes, por sua vez, acionam uma alma e um espírito, então, o velho vento
do sopro nômade que a diáspora carrega são as histórias e as palavras
esparsas, fagulhas e chispas de papéis. Territorialidades de trânsito
passíveis de apropriação e de pertencimentos postas em palavras. Textos
que são as chamas de vida e de celebrações judaicas religiosas e
diaspóricas. E que venham ventos mais leves.

Por fim, a \emph{Shechiná}, a presença e a habitação da divindade, que
deixou o Templo Sagrado (e que passou a habitar as palavras, os textos e
as interpretações), foi espalhada, peneirada, penetrada e permeada pelo
mundo e pela diáspora. Cabe a nós buscá-la, até que estes sentidos se
desvaneçam e sejam refeitos em outros territórios.

\begin{quote}
Os ventos ensinam, como diz um velho provérbio dos congos, que os
pássaros têm asas porque elas lhes foram passadas por outros
pássaros\footnote{\textsc{simas}; \textsc{rufino}, \emph{Encantamento. Sobre política de
  vida}, p.\,15.}.
\end{quote}

\begin{center}\adforn{49}\end{center}

\noindent{}Escrever é riscar as páginas. Este estilo riscante, como nos lembra
Vilém Flusser, se apresenta enquanto um estilete que dilacera e
esfarrapa imagens, ao final, ``as inscrições são imagens
esfarrapadas e rasgadas, imagens que o mortífero estilete
sacrificou''\footnote{\textsc{flusser}, \emph{op.\,cit.}, p.\,30.}. O estilete risca e
penetra a pedra ou o mármore para escrever, gravar nele suas breves
passagens, suas memórias, seus anseios e seus desejos de adeus e de
amor. No aramaico e no hebraico (ou no árabe), o movimento do estilete
sobre a pedra --- como nas \emph{Tábuas da Lei} --- é feito da esquerda
para a direita, em direção ao coração\footnote{Em uma tradução nunca se
  traduz palavras, mas sim sentimentos. Martin Buber nos conta que:
``O Rabi Mendel cuidava de que seus hassidim não usassem nada no
  pescoço durante a oração. Pois, segundo dizia, não deve haver
  interrupção entre o cérebro e o coração''. \textsc{buber}, \emph{op.\,cit.}, p.\,608.}.

A escrita, o ato de gravar e/\,ou depositar tinta sobre uma superfície,
também diz respeito à visualidade. Em ``Metaescrita'', Vilém Flusser usa
um feliz exemplo: ``\emph{palavra} é uma palavra, mas \emph{frase} não é
uma \emph{frase}. Só é possível reconhecer isso na escrita, pois uma sentença
como essa, se pronunciada, soa incoerente''\footnote{\textsc{flusser}, \emph{op.\,cit}., p.\,20.}.

Assim, ``escrever é afirmar o paradoxo de que, como no gesto
divino, as palavras criam o mundo mas não dão conta de aprendê-lo no
livro''\footnote{\textsc{manguel}, \emph{op.\,cit.}, p.\,33.}. Em \emph{O Golem de
Praga}, livro originalmente em iídiche do Prêmio Nobel de
Literatura (1978) Isaac Bashevis Singer, urge uma tentativa de criação
de mundos por meio da escrita e da palavra.

Em meados do século \textsc{xviii}, o \emph{Maharal} (Grão-Rabino) de Praga, Rabi
Loeb, confecciona, a partir do barro, uma matéria amorfa e sem alma ---
um \emph{golem} (um ``boneco'') ---, com o intuito de ajudar a
comunidade judaica local perseguida pelos \emph{pogroms}\footnote{\emph{Pogrom}
  (``massacre'') é uma palavra russa que se refere às perseguições
  deliberadas e aos atos em massa de violência, espontânea ou
  premeditada, contra as comunidades judaicas no Leste Europeu.}. Sobre
a testa da criatura, o Rabi inscreve a palavra hebraica \emph{emet} (em
português, ``Verdade'') e assim a figura amorfa --- proveniente da terra
--- ganha vida, ajudando a comunidade nas tarefas diárias. Mas o boneco
logo escapa às instruções do Rabi, obrigando o \emph{Maharal} a
devolvê-lo ao pó. Para isso, Rabi Loeb apaga a primeira letra da palavra
hebraica \emph{emet} ({אמת}), transformando-a em \emph{met}
({מת}), em português: ``morto''. O boneco então desmancha-se ---
retornando ao barro\footnote{Rabi de Guer costumava dizer que o homem
  foi posto no mundo para que transformasse o pó em espírito.}.

Dessa forma, ``embora a primeira letra da inscrição em sua testa
tenha sido apagada e \emph{emet} tenha virado \emph{met}, uma palavra continua ao
nosso dispor para nomear mais um inominável: a própria
morte''\footnote{\textsc{manguel}, \emph{op.\,cit.}, p.\,45.}. E não seria a morte o
objetivo de toda a criação? E não seria Deus ``aquele que renova
diariamente a obra da criação''\footnote{\textsc{buber}, \emph{op.\,cit.}, p.\,647.}?

A partir da presente narrativa, entende-se que, acima de tudo, o
artifício mágico da palavra cria, descria e recria mundos, assim como os
destrói --- para então renascer. A palavra cria o \emph{golem} --- ou o
ser humano, por meio da palavra divina ---, atesta a sua existência, o
nomeia, como também reitera a sua insignificância ao inscrevê-lo
novamente, ao destruí-lo.

A bênção judaica \emph{shehakol nihiyah bed'varo} (em português:
``por cuja palavra todas as coisas vieram a existir'' ou ``que tudo
tenha sido feito a Sua palavra'') é entoada sobre todos e quaisquer
alimentos (antes de seu consumo) que não provenham diretamente da terra,
por exemplo: peixes, carnes, frangos, ovos, queijos, mel e todas as
bebidas (exceto vinho) --- além de comidas consideradas mágicas, tais
quais os cogumelos e o chocolate.

Se Deus moldou o mundo --- e todas as coisas que ele contém --- por meio
do cruzamento entre a terra e a Sua mão (que confeccionou o barro),
então, todas as demais coisas que não são diretamente provenientes da
terra --- como os animais, embora vivam nela --- foram, necessariamente,
feitas à Sua palavra --- por meio do sopro.

Há algo curioso ao analisarmos o início do texto bíblico. Se partirmos
do pressuposto de que \emph{Berechit Bará Elohim et Hashamaim ve et
HaHaaretz}\footnote{\emph{Berechit Bará Elohim et Hashamaim ve et
  HaHaaretz} é a primeira frase da narrativa bíblica, traduzida para o
  português como: ``No princípio, Deus criou o céu e a Terra''.} e o
conectivo \emph{et} ({את}) entre as palavras \emph{Elohim} e
\emph{Hashamaim} seria desnecessário e está escrito com a primeira
({א}) e a última letra ({ת}) do alfabeto hebraico, então Deus
criou, antes de mais nada, tudo que está entre a primeira e a última
letra, e por fim, \emph{Hashamaim} --- o céu --- \emph{ve et HaHaaretz}
--- e a terra. Deus finalmente nomeou todas as coisas que existem, e
assim todas as coisas passaram a existir por conta de suas palavras e de
seus nomes, já que, segundo Alberto Manguel, ``é possível
articular em sílabas tudo que existe''\footnote{\textsc{manguel}, \emph{op.\,cit.}, p.
  35.}.

% Comentário da editora: Normalmente é escrito com "s", "Bereshit". Os irmãos Campos, no livro da Perspectiva, traduzem como "Bere'shith". Enfim, o hebraico do livro inteiro precisa ser revisto e padronizado (não há uma norma acordada, mas eu tenho a minha própria, com base no que tem sido feito). Vale também para a acentuação conforme fonética portuguesa.

\begin{center}\adforn{49}\end{center}

\noindent{}Certa vez, ``perguntaram ao Rabi Baruch: Por que dizemos: Bendito
o que falou, e o mundo se fez, e não Bendito o que criou o mundo?. Ele
respondeu: Louvamos a Deus porque criou nosso mundo com palavras, e
não, como todos os outros mundos, com o pensamento''\footnote{\textsc{buber},
  \emph{op.\,cit.}, p.\,131.}.

\begin{center}\adforn{49}\end{center}

\noindent{}Na história bíblica de Babel, a humanidade tentou chegar a Deus por meio
da construção de uma torre que alcançasse os céus. Manguel, a partir da
leitura de \emph{Gênesis}, nos conta que, neste cenário, Deus
``fragmentou o idioma único que toda a humanidade falava até o
momento na miríade de idiomas que falamos hoje''\footnote{\textsc{manguel}, op.
  cit., p.\,35.} --- o que inclui todas as línguas, vivas ou mortas. Em
contrapartida o Rabi Pinkhas sustenta que, ``antes da construção
da Torre, todos os povos tinham em comum a língua sagrada, mas além
dela, cada um tinha também o seu próprio idioma. Por isso é que está
escrito: \emph{Toda a terra tinha uma língua}, que era a sagrada, \emph{e alguns
falares}, que são as línguas adicionais e separadas de cada povo. Nestas
cada povo se entendia entre si, naquela os povos se entendiam entre si.
O que Deus fez, ao castigá-los, foi tirar a língua
sagrada''\footnote{\textsc{buber}, \emph{op.\,cit.}, p.\,177.}.

No conto ``O Congresso'', do escritor Jorge Luis Borges, ``um
homem sonha em compilar uma enciclopédia completa do mundo e no fim se
dá conta de que a enciclopédia em si já existe --- e é o próprio
mundo''\footnote{\textsc{manguel}, \emph{op.\,cit.}, p.\,43.}. Assim, como ocorreu na
história de \emph{Babel}, é impossível alcançar Deus porque Deus, em si,
é o próprio mundo.

Para a \emph{Cabalá}, Deus possui 72 nomes apresentados em três
versículos codificados. HaMakom e HaShem são apenas dois deles. HaShem,
em português, significa ``o Nome''. Já HaMakom significa ``o Lugar''.
Assim sendo, Deus é o Nome e o Lugar em que tudo ocorre, no qual existe
tudo o que existe. Deus abriu um espaço em si para a criação do mundo
--- e de todas as coisas que ele contém --- bem como da humanidade, ou
seja, ``a realidade de Deus e a representação que Deus faz da
realidade são idênticas''\footnote{\emph{Ibidem}.}.

Martin Buber nos conta que, certa vez, ``perguntaram ao Rabi
Pinkhas: Por que é que Deus é chamado de Lugar? Ele é certamente o
Lugar no Universo, mas deveriam então chamá-lo assim, e não apenas
Lugar? Ele responde: O homem deve entrar em Deus, de modo que Deus
possa envolvê-lo e tornar-se o seu Lugar''\footnote{\textsc{buber}, \emph{op.\,cit.}, p.\,167.}.

Se Deus é o Lugar em que tudo acontece e se materializa, então, a
escrita e as palavras materializam o texto no lugar. Pois ``há um
momento no qual descobrimos que das marcas de tinta na página emerge um
mundo completamente formado e magicamente real''\footnote{\textsc{manguel}, \emph{op.\,cit.}, p.\,34.}. É por meio do fino risco da caneta que a tinta, ao
chocar-se com o papel, transforma-se em lugar.

\begin{center}\adforn{49}\end{center}

\noindent{}Para a escritora e poeta chicana Gloria Anzaldúa (1942--2004) o termo
\emph{nepantla} diz respeito ao momento da escrita em que o todo é
caótico, ou melhor:

\begin{quote}
é um dos estágios da escrita, o estágio em que você tem todas essas
ideias, todas essas imagens, frases e parágrafos, e onde você tenta
transformá-los em uma única peça, uma história, enredo ou o que quer que
seja\footnote{\textsc{anzaldúa}, \emph{Borderlands/\,La Frontera}, p.\,237. Tradução da autora.}.
\end{quote}

\emph{Nepantla} é um espaço-tempo da construção da escrita, um termo que
se refere ao processo de estar entre meios. É um limiar em que múltiplas
formas de realidades (e de escritas) são vistas ao mesmo tempo ---
possivelmente enquanto um longínquo e interminável ensaio. Pois
``um ensaio é uma tentativa de incitar os outros a refletirem, de
levá-los a escrever complementos. (\ldots{}) Ele deve fazer rolar uma bola de
neve, na qual os complementos encubram cada vez mais a exposição
original''\footnote{\textsc{costa} \emph{apud} \textsc{flusser}, \emph{op.\,cit.}, p.\,11.}.

Segundo Gloria Anzaldúa\footnote{\textsc{anzaldúa}, \emph{op.\,cit.}, p.\,237. Tradução da autora.}, \emph{nepantla} é uma palavra do universo indígena
\emph{nahuatl} e indica algo que está entre dois corpos d'água --- um
espaço-tempo entre dois mundos. Este território é um limiar, um local
que não está nem aqui, nem lá, mas em transição.

Se as coisas são escritas entre meios, então \emph{nepantla} é um
espaço-tempo de fronteira, em que continuamente (e constantemente) se
encontra o mundo. São espaços-territórios --- geográficos ou
linguísticos --- de encontro com a alteridade. Assim, habitar o entre
resulta numa visão dupla, primeiro da perspectiva de uma cultura, depois
da perspectiva de outra. Ao viver entre fronteiras, as comunidades
judaicas, por exemplo, bem como os demais grupos socialmente
minoritários (chicanos, quilombolas, indígenas, refugiados etc.), acabam
vendo duplamente, interagindo com estes ambientes tanto no âmbito do
texto, quanto no âmbito do espaço, e criando locais de pertencimento.

\begin{center}\adforn{49}\end{center}

\noindent{}Na tradição da escrita judaica religiosa --- e de suas interpretações
---, as páginas, os papéis e os pergaminhos --- bem como os seus textos
--- são territorialidades de apropriação e espaços sagrados. Este texto
escrito e inscrito sobre um pedaço de papel, pano, couro ou celulose
influencia as comunidades, suas relações com a divindade e com os
mandamentos divinos. Nos mundos das narrativas e dos livros, a
manifestação do texto no espaço\footnote{As mais diversas interpretações
  dos textos sagrados judaicos se manifestam no espaço, pois conduzem as
  comunidades judaicas na diáspora e fora dela e acarretam mudanças
  reais nas maneiras de manifestar a judeidade, bem como nas dinâmicas
  dos territórios em que as comunidades se agregam.} --- como é o caso
da leitura dos rituais e das festividades judaicas religiosas ---
transforma as comunidades, seus espaços de agregação, bem como as suas
histórias. Em uma perspectiva de fronteira linguística, a ficção dos
textos, segundo o pesquisador e filósofo Leonel Olímpio:

\begin{quote}
Também implica no problema de como aquele mundo ficcional nos
influencia, e também até que ponto nós influenciamos aquele mundo
ficcional. Como assim? Por exemplo, se partirmos do pressuposto de que o
mundo daqui (presumidamente ``real'') não influencia o outro, toda
leitura, toda linguagem exposta em qualquer texto não mudaria de
interpretação com o tempo\footnote{OLIMPIO, \emph{Há o Anti-Limite?
  Sobre ficções, equívocos e metafísica experimental}, p.\,67.}.
\end{quote}

Neste sentido, os textos judaicos sagrados são passíveis de
interpretações, e o leitor é um participante ativo das páginas --- o que
faz com que os escritos se tornem territórios fluidos, espaços de
diálogo e de apropriação. O leitor de textos judaicos não lê
simplesmente palavras escritas em uma página; o leitor é um ator ativo e
um participante vocal na conversa.

O debate das escrituras ocorre em meio às páginas, e seus assuntos
abrangem séculos, milênios e gerações de professores, mestres, rabinos e
pensadores talmúdicos. O Talmud (em português: ``aprendizado''),
livro composto pela \emph{Mishná} (``estudo repetido'') e pela
\emph{Guemará} (``conclusão''), é uma série de textos-ensaios que
procuram --- a todo custo --- desmembrar as histórias (\emph{reductio ad
absurdum}), as regras, as restrições e as leis judaicas religiosas a
ponto de dilacerar suas minúcias.

Os textos sagrados, bem como as interpretações rabínicas, são escritos
às vezes em hebraico e às vezes em aramaico, sem pontuações. Nunca há
página de número 1 em um texto judaico tradicional, os textos sempre
começam na página 2, já que no Talmud não existe nem o começo
nem o fim, tudo é uma questão de meios. O texto principal é posto ao
centro da página rodeado de comentários dos mais diversos estudiosos ---
de diferentes territórios, gerações e contextos. A diagramação é feita
de tal maneira que não há hierarquia entre as interpretações rabínicas,
as exposições dos argumentos e os conceitos apresentados, ou seja, todos
os estudiosos e estudiosas são convidados a trazer uma nova
interpretação das leis judaicas.

Alguns comentadores escreveram tratados, outros apenas alguns rabiscos
sobre uma folha de rascunho, mas no Talmud tudo serve para
tornar as discussões mais dinâmicas. Dessa forma, as mais diversas
visões rabínicas são apresentadas enquanto continuidades dos próprios
textos sagrados, abrangendo pelo menos dois mil anos. Por fim, as
perspectivas rabínicas sobre os textos religiosos e sagrados --- suas
interpretações --- são condicionadas por suas épocas e momentos
históricos, suas relevâncias para a comunidade e seus sentidos judaicos.

\begin{center}\adforn{49}\end{center}

\noindent{}Certa vez, ``perguntaram ao Rabi Mihal: Está escrito na Ética dos
Pais: Quem é sábio? Aquele que aprende com todos os homens, como está
escrito: De todos os meus mestres recebi entendimento. E por que então
não está escrito: Aquele que aprende com todos os mestres?. O Rabi
Mihal explica: O mestre que pronunciou tais palavras pretende explicar
que se deve aprender não só com aqueles que atuam como mestres, mas com
todos os homens. Mesmo com o ignorante, mesmo com o perverso, te é dado
alcançar o entendimento de como conduzir a vida''\footnote{\textsc{buber},
  \emph{op.\,cit.}, p.\,189.}.

Os estudos sobre os livros e textos judaicos religiosos são feitos em
parceria --- \emph{chavruta}\footnote{\emph{Chavruta} é uma palavra
  hebraica proveniente da raíz \emph{chaver}, em português
  ``amigo''. E é uma prática de estudo em que os textos são lidos em
  pares.}, pois se faz necessário falar a respeito do texto como se ele
estivesse realmente vivo. Ao contrário das bibliotecas modernas, lugares
de culto ao silêncio, o \emph{Beit Midrash} (em português: ``Casa de
Estudos'') é o local onde os debates são acalorados, ruidosos, animados
e cheios de gesticulação e de frustrações. A Casa de Estudos reflete o
livro --- o tumulto e o barulho de cada uma das interpretações.

Entende-se aqui que há uma incapacidade de aprender o mundo por completo
em uma só palavra (ou em uma só opinião), pois ``um poema nunca é
concluído, apenas abandonado''\footnote{\textsc{manguel}, \emph{op.\,cit.}, p.\,42.}.
Como observa Flaubert, ``a estupidez consiste em desejar
concluir''\footnote{FLAUBERT \emph{apud} ibidem, p.\,46.}. Ou seja, o texto
será tanto mais significativo quanto maior o número de leituras e de
interpretações. E assim, nem mesmo as diversas interpretações de
\emph{Bereshit} chegam a uma conclusão. Muitas são as traduções e
interpretações de ``No princípio era o verbo'', porém, o mais
correto seria ``Deus começa a criar'', na medida em que o
judaísmo não produz construções que possam levar à idolatria. Rejeita, a
todo custo, a tentação da serpente em aspirar à condição divina, de
fixar e de sacralizar conceitos em uma só conclusão. Se partirmos do
pressuposto de que ``No princípio era o verbo'', então
automaticamente estamos colocando um começo, fixando imagens. Como
resposta todos os livros e, em especial, a máxima literária judaica, a
Torá, são passíveis de interpretações e de apropriações.

% Comentário da editora: De onde é "no princípio era o verbo"? Do próprio Bereshit?

Uma história talmúdica nos conta que, em meio a uma discussão
rabínica, dois indivíduos, buscando defender seus pontos de vista,
iniciam um debate a respeito de uma lei judaica. O primeiro rabino diz:
``Se eu estiver correto que esta árvore o prove''. Depois de
alguns segundos, passa um vento sobre a árvore fazendo-a balançar. O
segundo rabino responde: ``Não há como comprovar que essa árvore
balançou por conta de sua suposição''. O primeiro rabino novamente diz:
``Se eu estiver correto que esta pedra o prove''. Logo em
seguida, a pedra desliza penhasco abaixo em direção ao mar. E novamente
o segundo rabino responde: ``Não há como comprovar que a pedra
deslizou por conta de sua suposição''. O primeiro rabino, impaciente,
diz: ``Se eu estiver correto que os céus o provem''. Nesse
instante a esfera divina brilha, permitindo uma passagem de luz. O
segundo rabino, deslumbrado, observa e diz: ``O poder já não está
mais nas mãos dos céus. A Torá foi dada por Deus aos seres humanos e
agora está sob a vontade deles!''\footnote{\textsc{oz; oz-salzberger}, \emph{Os
  judeus e as palavras}, p.\,18.}.

Durante os últimos dois mil anos, e ainda hoje, as discussões e os
diálogos manifestaram-se através do livro judaico; o debate judaico não
ocorre em fóruns abertos, mas nas páginas escritas. Seus assuntos
abrangem séculos, até mesmo milênios, e convergem de maneiras
excepcionais, criam espaços, territórios de apropriação.

\begin{center}\adforn{49}\end{center}

\noindent{}O \emph{Midrash}\footnote{O \emph{midrash} é uma metodologia judaica de
  leitura das escrituras sagradas. Seu fundamento está em uma visão
  cíclica da história, segundo a qual o acontecimento central, ou
  melhor, a presença de Deus no texto, é continuamente revisado,
  oferecendo sempre novos significados às várias situações expostas, bem
  como às interpretações.} diz que toda e cada letra da Torá deve
estar completamente envolvida por um pedaço de pergaminho\footnote{Referência
  à \emph{Menachot} 29a}. O pergaminho branco em volta das letras pretas
é parte integrante dos escritos e do rolo da Torá. Em debates
talmúdicos extremamente aprofundados, os sábios interpretam que a
Torá é ``escrita em fogo negro, sobre fogo branco''. As
separações e os espaços brancos do pergaminho permitiram que Moisés,
durante a inscrição da fala de Deus no texto, analisasse e absorvesse os
aprendizados anteriores, para assim seguir ao próximo. O fogo branco é o
espaço mais elevado da contemplação, ao passo que o fogo negro (a
escrita) é a revelação do intelecto no domínio da linguagem.

\begin{center}\adforn{49}\end{center}

\noindent{}A cultura e a filosofia judaicas, bem como as suas tradições, conhecem
muito bem a relevância do espaço da página, dos textos e dos livros para
a construção de territorialidades. A divindade judaica habita o texto,
assim como também habita a recordação --- a lembrança\footnote{Por mais
  que este trabalho se sustente em uma perspectiva judaico-diaspórica
  sobre o lugar, ou seja, em territorialidades judaicas móveis,
  transitórias e fluidas --- sendo elas: \emph{BS''D}, \emph{kipá},
  \emph{mezuzá}, \emph{eruv} --- há um território espiritual e simbólico
  que permeia todas as discussões: a memória e a lembrança do Templo
  Sagrado de Jerusalém. Ao final, todas as rezas, orações, tradições e
  cultos judaicos religiosos são conduzidos à cidade de Jerusalém, local
  sacrossanto judaico --- território da construção do Templo ---, e
  assim levados ao céu.} --- e os anseios de retorno do Templo Sagrado
de Jerusalém.

As territorialidades sagradas dos textos e das páginas judaicas
religiosas são cosmogonizadas por meio da inscrição ao canto da folha do
artifício de {בס''ד} (BS''D). Em aramaico, a inscrição BS''D
(acrônimo de: \emph{B'Syata D'bishmaya}) significa ``com a ajuda
dos céus''. Em hebraico moderno, a mesma inscrição significa: ``em uma
cinta''. Este artifício passou a ser introduzido pelo Rabi Yehuda
HaChassid entre os anos de 1150 e 1217, enquanto um símbolo de lembrança
divina. Ou seja, quando realizamos tarefas cotidianas e mundanas Deus
deve estar em nossas mentes, e em nossas penas e canetas\footnote{A
  inscrição da tinta de uma caneta sobre uma página cria um território.}.

O judaísmo entende que as páginas --- estes espaços em branco, prontos
para serem riscados --- são locais de contemplação e de criação. Ao
realizar um traço, um desenho e/\,ou um rascunho sobre um pedaço de papel,
os seres humanos criam um mundo e um universo, por completo --- seja ele
virtual, ficcional ou real. Escrever é criar mundos, e o mesmo pode ser
dito de desenhar, gravar, inserir ou apagar, como nos lembra a história
do \emph{golem}.

Assim, há o costume de gravar o acrônimo BS''D nas páginas,
anteriormente a quaisquer inscrições. Este artifício relembra que apenas
``com a ajuda dos céus'' é possível criar --- ao final, nada está acima
da criação divina, nem mesmo a fruta comida por Adão no Jardim do
\emph{Éden}.

Por fim, ambas as traduções --- ``com a ajuda dos céus'' (do
aramaico ao português) e ``em uma cinta'' (do hebraico moderno ao
português) --- relacionam-se com a Torá. A Torá é um
material escrito que dialoga com a existência, com a crença e com a
presença de Deus (por meio do texto) no mundo humano. Este
livro-documento é envolto e fechado por uma cinta --- um pedaço de pano
ou de couro --- e encapsulado em uma caixa de madeira. Assim, a tradição
judaica se estabelece no enlaçamento (cinta) entre a presença divina, a
fidelidade e a lealdade judaica ao monoteísmo (com a ajuda de Deus) e a
fé no livro.

\begin{center}\adforn{49}\end{center}

\noindent{}O \emph{aleph} ({א}) é a primeira letra do alfabeto hebraico --- e do
alfabeto aramaico ---, usado para escrever o iídiche e o ladino.
É uma letra que contém dentro de si todas as demais letras --- contém o
universo inteiro, sem fim, e ``tem a forma de um homem que
assinala o céu e a terra, para indicar que o mundo inferior é o espelho
e o mapa do superior''\footnote{\textsc{kilsztajn}, \emph{Yiddish}, p.\,20.}. O
\emph{aleph} ({א}) é o Lugar onde estão, sem se confundirem, todos os
lugares do mundo, vistos por todos os ângulos --- onde o interlocutor vê
o que está nascendo e partindo.

Certa vez li que a palavra \emph{vayikra} --- {ויקרא} (em português:
``Ele chamou'') --- está escrita de uma forma pouco comum na
Torá. A última letra, o \emph{aleph}, está em miniatura. Os
sábios interpretam que a tentativa de esconder esta letra foi uma
expressão de extrema humildade e de devoção de Moisés a Deus. Por assim
dizer, Deus apenas ``aconteceu'' (\emph{vayikar}) --- {ויקר}.

\begin{center}\adforn{49}\end{center}

\noindent{}O ato de escrever, a incisão feita ao gravar letras, palavras e frases
sobre uma página em branco, cria mundos, realidades e universos por
completo. Este incansável e interminável movimento é feito neste exato
instante e neste exato documento\footnote{Pretende-se que as diversas
  interpretações dos textos, aqui expostos, ressoem em você que lê, o
  leitor.}. E assim, como intérprete de uma longa tradição de textos e
de livros, acredito que o nosso mundo --- e todas as coisas que ele
contém --- está sendo criado e reinventado por nós a todo instante.

Por fim, os textos aqui escritos e inscritos --- suas notas de
rodapé\footnote{A disposição gráfica das notas de rodapé alude à página
  dos textos judaicos religiosos, bem como do pensamento rabínico que há
  no Talmud. Pois, no Talmud, a menção de um nome,
  situação ou ideia permite a introdução de uma história que alivia o
  clima de discussões complexas e leva o debate adiante. As notas, ou
  melhor, os comentários e as interpretações rabínicas, rondam,
  circunscrevem e flutuam ao redor do texto principal sem
  hierarquização. Aludindo à organização talmúdica, procuro
  incitar os outros a refletir, abrindo espaço para as mais diversas
  interpretações, histórias e vozes.} e suas interpretações
(ressonâncias) que levam a outras histórias, anedotas e contextos ---
atestam a veracidade da fala de Rabi Tarfon: ``Não lhe é exigido
que complete a tarefa, mas não é livre para escapar dela''\footnote{\emph{Pirkei
  Avot} 2:16, escrito cerca do século \textsc{i}.}.

A diagramação destas páginas-territórios são releituras dos textos
sagrados, uma tentativa de apropriação das chispas e das fagulhas que
voam e contornam o território do texto --- já que o espaço do texto é um
território. E mais que nada, neste instante, é hora de nos apropriarmos
dele.

\begin{center}\adforn{49}\end{center}

\noindent{}``No mundo da aprendizagem, das discussões e dos debates judaicos,
as palavras dos textos, literal e figurativamente,
ressoam''\footnote{HALPERN, \emph{The Art of Debate: Jewish Style}. Tradução da autora.}. E, ao ressoar, entoam uma voz --- um sopro de
espírito, uma alma. As palavras, finalmente, cosmogonizam o
mundo\footnote{As vozes não se calam, os sopros fazem eco para a
  eternidade.}.

Desde outubro de 2023 \emph{ya no ay biervos}\footnote{Do ladino: ``já
  não há palavras''.}. ``No estupendo Nossa Música, de Jean-Luc
Godard, uma jornalista israelense entrevista o poeta palestino Darwish,
que, privado de sua terra, fez das palavras a sua pátria. E ela comenta:
\emph{você começa a soar como judeu!}. O devir-judeu do palestino, o
devir-palestino do judeu''\footnote{\textsc{pelbart}, Peter Pál. \emph{Ser judeu
  no Brasil}.}.

Por fim, todos nós, antes de mais nada, vivemos em palavras.

\chapter{Corpos}

\noindent{}O filósofo e historiador Martin Buber conta que certa vez perguntaram ao
Rabi Iehiel Mihal\footnote{Rabi Iehiel Mihal era conhecido como
  \emph{Tzadik HaGamur} (um Justo) e foi um dos mais notáveis discípulos
  do Rabi Baal Shem Tov, precursor do movimento \emph{hassídico} judaico
  no Leste Europeu.} a razão de retardar o começo de suas orações. O
rabino respondeu: ``Como a tribo de Dan\footnote{Após o
  recebimento das Tábuas da Lei, o povo hebreu foi dividido em doze
  tribos, sendo cada uma delas responsável por um pedaço de terra em
  \emph{Canaã} (antigo território que hoje engloba Israel, Palestina,
  Líbano, Cisjordânia, Jordânia e o Deserto do Sinai). As doze tribos
  são: Reuven, Shimon, Levi, Yehuda, Dan, Naftali, Gad, Asher,
  Issachar, Zebulun, Yossef e Biniamyn.}, que na Terra de Canaã
caminhava atrás de todas as demais tribos, recolhendo tudo aquilo que
perdiam; assim também sou eu --- um coletor incansável de migalhas de
nossa, e minha, própria história''\footnote{\textsc{buber}, \emph{op.\,cit.}, p.\,192.}.

Tal qual a tribo de Dan, Rabi Iehiel Mihal e Martin Buber, no judaísmo
todos os indivíduos são convidados a fazer um \emph{chidush}\footnote{Na
  literatura rabínica, o \emph{chidush} refere-se a uma nova abordagem
  de ideias e interpretações em obras já existentes. É uma forma de
  inovação feita dentro do sistema da \emph{Halachá} --- série de
  documentos que descrevem as maneiras judaicas de conduzir a vida, suas
  regras e verdades.} e podem, assim, se apropriar e interpretar os
textos sagrados e bíblicos\footnote{Na tradição judaica todos os
  indivíduos possuem o mesmo nível de conexão com a divindade, sem a
  necessidade de intermediação.}. Para se tornar este coletor de
migalhas --- vestígios de uma história particular compartilhada ---,
como nos lembram o escritor Amós Oz e a historiadora Fania
Oz-Salzberger, não é preciso ser filha ou filho de um ventre judaico,
basta ser um leitor --- um apropriador incansável, um intenso
talmúdico, um intérprete.

Das mais dispersas migalhas de narrativas, individuais e coletivas,
agarradas em meio ao trânsito, essas múltiplas referências, por fim, se
tornam sementes e são germinadas em determinadas territorialidades,
agregando influências e interferências dos locais que passam a
frequentar. Pois, no campo dos Estudos Judaicos, entende-se que
historiadoras e historiadores não produzem uma História Judaica por
excelência, mas sim histórias --- narradas e escritas --- sobre as
pegadas destes judeus e judias nos territórios que percorrem.

Neste sentido, este texto sustenta-se enquanto um novo começo, de outras
--- mas nunca novas --- interpretações passíveis de apropriação e de
pertencimento. Como todas as evidências são escoradas por mistérios,
entende-se aqui que se não se fizer o fogo não há o que ver, ou melhor,
não há de se ver nada que se deseja. Porém, ao se decidir por acendê-lo
então \emph{kudiado con hacer fuego, i no aguentar el kalor}.

\begin{center}\adforn{49}\end{center}

\noindent{}Em ``Tempo e atemporalidade'', terceiro capítulo do livro \emph{Os
judeus e as palavras}, algumas dessas ``clareiras-fogos'' já se mostram
acesas. Amós Oz e Fania Oz-Salzberger discorrem a respeito do conceito
de espaço e de tempo para a civilização, para a tradição e para a
cultura judaicas. Parte-se do pressuposto de que todas as civilizações
estão preocupadas com o seu passado e no judaísmo isso não poderia ser
diferente. São as orações e as rezas, os cultos, os ritos, as tradições,
a cultura e as festividades que lembram e relembram cotidianamente o
contínuo \emph{quórum} --- em que a prática levaria à mística\footnote{Os
  ritos judaicos baseiam-se em uma série de preceitos e de tarefas
  religiosas realizadas cotidianamente. A \emph{kavaná} --- a elevação
  mental e espiritual que o indivíduo atinge ao executar os deveres
  religiosos --- está também entremeada de uma rotina, em que a
  constante prática e o refazimento são essenciais para alcançar a
  santidade.} --- de uma tradição ancorada em textos e em livros.

Assim, em parte, é o fio genealógico --- gênico ou textual --- que
transmite as mais variadas interpretações e tradições judaicas, de
geração em geração --- e firmam-se, finalmente, para que então possam
ser lidas e tateadas em outros territórios e em outros
tempos\footnote{Aqui vale lembrar que os documentos, as tradições e as
  comunidades, ao firmarem-se em outras territorialidades e tempos,
  ainda são passíveis de indeterminabilidade, condicionados aos
  movimentos das diásporas voluntárias ou forçadas. Ao serem levados
  novamente pelo velho vento do sopro nômade da diáspora, brotarão em
  outros locais, em outros tempos, em um movimento interminável, cíclico
  e constante.}. Em ressonância, todos os dias da semana, ao colocar o
\emph{tefilin}\footnote{O \emph{tefilin} é um objeto de cunho religioso
  que consiste em um par de caixas unidas por meio de uma tira de couro
  animal. Conhecido como filactério, o \emph{tefilin} é um instrumento
  de proteção e fortificação da aliança divina usado por homens, e mais
  recentemente incorporado por mulheres, nos momentos de culto e rito
  religioso. Esse objeto é utilizado apenas durante os dias de semana,
  não sendo incorporado nas práticas judaicas do \emph{shabat}.}, judeus
e judias selam a devoção e afeição entre Deus e o povo judeu. Na
Torá, Deus e povo são aludidos enquanto um casal, um par --- e,
em união, simbolizam a aliança, a lealdade intelectual e uma
espiritualidade e religiosidade atreladas ao divino\footnote{Segundo o
  Talmud, Deus criou os seres humanos --- e todas as coisas do
  mundo --- para terem prazer com a conexão divina. Todas as pessoas são
  passíveis de manifestar o divino em suas vidas por meio de ações
  ordinárias e extraordinárias. Pois entende-se que o divino --- para
  além de Deus --- é a presença onipotente e onisciente, e é ele,
  propriamente, a matéria.}.

É a partir do conceito de que a divindade está permeada de livros, de
textos e de ritos --- em especial uma lealdade intelectual --- que o
Rabi Akiva\footnote{Rabi Akiva (50 d.E.C.--135 d.E.C.) foi um renomado
  estudioso hebreu e introduziu um novo método de interpretação das leis
  e das tradições orais judaicas, que veio a se tornar a \emph{Mishná}. A \emph{Mishná} é uma das principais obras do judaísmo rabínico
  e a primeira grande redação da tradição oral judaica.}, ao se deparar
com a destruição do segundo Templo de Jerusalém\footnote{O segundo
  Templo Sagrado de Jerusalém, ou Templo de Jerusalém, foi uma série de
  construções de cunho religioso que se localizava no Monte
Moriá (cidade velha de Jerusalém), atual local do Domo da
  Rocha e da Mesquita \emph{Al-Aqsa}.} (76 d.E.C.), dança incansável e
insaciadamente em torno de seus escombros e resíduos. É instigante
pensar que, das duas vezes\footnote{O Templo de Jerusalém foi destruído
  duas vezes, no ano de 587 a.E.C. e posteriormente no ano de 76 d.E.C.}
que se tentou --- talvez a todo o custo --- fixar a
\emph{Shechiná}\footnote{Conta o Talmud que, quando os judeus
  pecaram, a presença de Deus deixou o Templo Sagrado de Jerusalém e
  retirou-se em direção aos céus. Porém, uma parte da \emph{Shechiná}
  permaneceu em terra para acompanhar os judeus e judias que
  forçosamente deixaram a Terra de \emph{Canaã}, levados ao exílio na
  Babilônia. Conta a tradição que a presença divina --- não mais
  presidida em uma construção de pedra --- passou a ser transportada
  pelas crianças¹ em direção ao exílio. Outro exemplo diz: ``Está
  escrito na Guemará: Todas as datas previstas já passaram. Mas assim
  como a Shechiná deixou o Santuário e foi ao exílio no decurso de dez
  viagens, assim também ela não voltará de uma só vez e a luz da
  salvação paira entre o céu e a terra''. A \emph{Shechiná} foi levada ao exílio através das crianças,
  que, na tradição judaica, são instrumentos visíveis de esperança, de
  que haverá uma continuidade deste povo. \textsc{buber}, \emph{op.\,cit.}, p.\,386.}, a habitação e morada de Deus foi
destruída. Assim, a divindade que sobrevive a esta destruição --- e que
neste texto busca-se reconstituir --- está diretamente vinculada a um
conceito de Deus e de divino amparado por arquiteturas não mais estáveis
e intransponíveis, mas sim móveis, reprogramadas nos espaços e nos
tempos, apropriadas por cada uma das gerações judaicas e levadas em
bolsos, maletas, sacolas nas diásporas.

Ainda nesta chave, como nos lembram Amós Oz e Fania Oz-Salzberger:
Não existe \emph{terra santa} na Torá. O único lugar sacrossanto é o
Templo de Jerusalém''\footnote{\textsc{oz; oz-salzberger}, \emph{op.\,cit.}, p.\,121.}
--- motivo pelo qual todas as sinagogas estão voltadas nessa direção. A
partir das constantes destruições --- de arquiteturas e de comunidades
---, o renascer das cinzas das ruínas faz parte das vidas judaicas em
diásporas e se estabelece enquanto pedras --- estilhaçadas e
fragmentadas --- com tanta sede de celebração que são capazes de
inventar a vida só para poder beber dessa existência. 

E se essas pedras falassem, o que elas diriam?\footnote{Elas diriam que:
``No verdadeiro amor tudo é para sempre vivo. E sei que, como as
  pedras, existo pela sede. Quero sempre inventar a vida''. \textsc{mãe}, \emph{As mais belas coisas do mundo}, p.\,43.}) Antes mesmo
de selar um contrato matrimonial --- um dos principais motivos de
celebração das comunidades judaicas ortodoxas --- há o costume de o
noivo quebrar um copo de vidro contra o chão. Este ato é uma recordação
da destruição do Templo Sagrado e de que nunca se estará completo sem a
presença da \emph{Shechiná}\footnote{Conta Martin Buber: ``Quando
  o Rabi Baruch chegou, no salmo, às palavras: \emph{Não deixarei que durmam
  meus olhos, nem que descansem minhas pálpebras antes de ter encontrado
  uma morada para Deus}, parou e disse de si para consigo: \emph{Antes de eu
  me ter encontrado e transformado em morada pronta para o pouso da
  Shechiná}''. \textsc{buber}, \emph{op.\,cit.}, p.\,131.}. Por fim, quebra-se o copo e dança-se
em seus escombros, como fez o Rabi Akiva.

O judaísmo é feito de tempos santos. Eles se consolidam enquanto uma
espécie de calendário cíclico e elipsoidal que foge do binômio
progresso-tempo, abruptamente implementado pelo renascimento e pela
modernidade ocidental. Entre estes tempos santos estão o \emph{shabat},
os \emph{chaguim}, os jejuns, os dias de \emph{Yom Tov}, as mais
diversas e variadas celebrações e festividades cíclicas judaicas --- que
recordam desastres e perseguições, momentos de glórias e triunfos,
salvamentos divinos e outros salvamentos essencialmente humanos. Ao fim,
entende-se que, embora o Templo de Jerusalém tenha sido destruído, os
judeus ``ainda puderam levar consigo para a amarga diáspora suas
santidades não espaciais, intangíveis: a língua, as leituras e o sempre
recorrente, ciclicamente confortante calendário de
\emph{tempos-santos}''\footnote{\textsc{oz; oz-salzberger}, \emph{op.\,cit.}, p.\,121.}.

Estes tempos santos --- alguns recordados semanalmente, outros
anualmente\footnote{Alguns tempos sagrados judaicos são lembrados e
  realizados semanalmente, como o \emph{shabat} --- recordação da
  criação divina do mundo e de todas as coisas que ele contém. Outros
  eventos são recordados anualmente, como as festividades e os jejuns.
  Há ainda as recordações geracionais, que ocorrem a cada sete anos,
  como é o caso dos anos de \emph{shemitá}, em que há a proibição do
  trabalho agrário, permitindo-se assim apenas a coleta dos frutos e
  frutas.} --- são celebrações de rito e de vida partilhadas e
realizadas em comunidade. Ao final, entende-se fortemente que os povos
em diáspora celebram a vida. E, como uma espécie de calendário,
espera-se que, um dia, no lugar de memórias difíceis e tristeza --- tais
quais os eventos de perseguição, de intolerância, de destruição ---
haverá celebração. Mas não esperaremos calados.

O artifício tão recorrentemente usado de celebrar a vida\footnote{O
  humor é uma característica intrínseca aos judeus \emph{ashkenazi}
  (judeus provenientes da Europa Central e do Leste Europeu). Apesar de
  ser uma particularidade relativamente recente e incorporada à cultura
  judaica no final do século \textsc{xix}, este atributo foi introduzido na
  literatura iídiche por meio do escritor Sholem Aleichem ---
  que, ao retratar o sofrimento do povo judeu, se viu forçado a manter
  acesa a chama da vida, da alegria e do humor de viver. Assim:
  ``{[}\ldots{}{]} Sholem Aleichem capturou a alma de um povo que se
  viu identificado e que sorriu em retorno para ele e aprendeu a sorrir
  com ele''. \textsc{kilsztajn}, \emph{op.\,cit.}, p.\,79.} --- um \emph{modus operandi}
judaico de ser e estar no mundo --- está presente nos festejos ortodoxos
e tradicionais, nas banalidades do cotidiano, na quebra do copo de vidro
ao selar um contrato matrimonial. Está firmado no sempre alto, e sempre
belo, \emph{Lechaim}\footnote{A palavra hebraica \emph{Lechaim}, usada
  nos momentos de celebração, é escrita e falada no plural (todas as
  palavras hebraicas terminadas com sufixos ``im'' ou ``ot'' são
  palavras no plural). Pois, no judaísmo, entende-se que o
  \emph{Lechaim} (no português: ``Saúde!'') deve sempre ser entoado
  coletivamente, em comunhão. Uma segunda interpretação diz que, quando
  celebra-se a vida, não se deve celebrar somente uma única vida. Todas
  as vidas são passíveis de celebração e de cerimônia. Dessa maneira, a
  língua hebraica compreende que o \emph{Lechaim}, melhor traduzido para
  ``Às vidas!'', para além de uma celebração compartilhada entre pessoas
  que situam-se no mesmo espaço-tempo, é uma forma de festejar e lembrar
  todas as demais existências criadas e moldadas por Deus.}. Como também
nas orações, nas bênçãos e nos cânticos entoados coletivamente --- em
comunhão --- nas mais dispersas e variadas sinagogas e comunidades. O
celebrar está nos ritos, na passagem de morte. São temporadas curtas,
mas alegres, que se repetem anualmente, e que são guardadas e
transmitidas pelos movimentos cíclicos e elipsoidais da continuidade de
uma tradição, e de uma educação judaica, que garante constantes
manutenções --- que atravessa escalas, e é ao mesmo tempo transversal a
elas --- de vínculos pessoais de identidade, que consolidam comunidades
nos mais diversificados territórios.

Dessa maneira, ao traçar este raciocínio, o \emph{chidush}, a
continuidade textual, a lealdade intelectual, o senso de comunidade, de
agregação, de coletivo, as leituras e os estudos bíblicos e
talmúdicos são preceitos que garantem vínculos identitários e
comunitários entre as diversas gerações que transitam pelo mesmo
espaço-tempo dentro das comunidades. Assim, esses sujeitos --- ao se
debruçarem sobre os textos sagrados --- transcendem as barreiras
temporais e criam para si outros mundos --- outros lugares de diálogo e
de estudos por meio das interpretações.

\begin{center}\adforn{49}\end{center}

\noindent{}\begin{quote}
``Conta o Baal Schem Tov\footnote{Rabi Israel ben Eliezer (1698--1760),
  também conhecido como Baal Shem Tov, foi um dos principais rabinos
  poloneses e está entre os precursores do judaísmo \emph{hassídico}.
  Baal Shem Tov significa em português: ``Um com um bom nome''. O judaísmo \emph{hassídico} é um movimento da mística judaica
  que dá ênfase à alegria religiosa e é organizado em grupos chamados de
  dinastias \emph{hassídicas}, governados por rabinos cuja liderança é
  herdada nas famílias.}. \emph{Dizemos: o Deus de Abraão, o Deus de Isaac e
o Deus de Yaacov, e não o Deus de Abraão, Isaac e Yaacov; pois Isaac e
Yaacov não se apoiavam nos estudos e serviços de Abraão, mas eles mesmos
buscavam a unidade do Criador e o seu serviço}''\footnote{\textsc{buber}, op.
  cit., p.\,95.}.
\end{quote}

A filosofia judaica parte do pressuposto de que, dentro de uma sala de
aula ou dentro de uma mesma família, todos os indivíduos --- sejam eles
estudiosos, historiadores, \emph{tzadikim}\footnote{\emph{Tzadik} (no
  plural: \emph{Tzadikim}) é um título concedido a personalidades do
  judaísmo ortodoxo consideradas santos, mestres e/\,ou líderes religiosos
  e espirituais.}, alunos, alunas e mestres --- são postos em debate e
convidados a trazer um novo contexto, novas interpretações e ideias
sobre os textos sagrados\footnote{Conta Martin Buber: ``Um
  discípulo perguntou ao Maguid de Zlotzchov: --- Está escrito no livro
  de Elias: \emph{Todos em Israel têm o dever de dizer: Quando chegará a
  minha obra aos pés da obra de meus patriarcas Abraão, Isaac e
  Yaacov?}. Como se deve entender isso? Como podemos atrever-nos a
  pensar que seríamos capazes de imitar os patriarcas? Explicou o Rabi:
  \emph{Assim como os patriarcas inventaram novas maneiras de servir, cada
  qual um novo serviço conforme sua peculiaridade, um pela mercê, o
  outro pela força, o terceiro pela glória, assim também nós devemos,
  cada um a seu próprio modo e à luz dos ensinamentos e do serviço,
  inventar inovações e não fazer o que já foi feito, mas o que ainda
  está por fazer}.'' \textsc{buber}, \emph{op.\,cit.}, p.\,190.}. Ao colocar estes diversos indivíduos
em um mesmo espaço --- seja ele físico, como uma sala de aula, ou
temporal, que transcende a linha cronológica --- a noção
temporal-espacial judaica se vincula diretamente às argumentações e
constatações dos documentos e das palavras ali escritas.

Nesta lógica, para além de espaços-tempos retilíneos, cronológicos e
progressivos, a noção judaica de espaço-tempo está relacionada
diretamente aos textos --- e aos livros --- e às suas perenidades de
interpretações. São convocadas a uma sala de estudos --- o
\emph{cheder}\footnote{O \emph{cheder} ocupa um papel central no
  folclore judaico, em especial no mundo iídiche, e simboliza a
  escola primária --- local de excelência da educação judaica, em que
  eram ensinados fundamentos do judaísmo e do hebraico bíblico. Os
  estudos eram realizados normalmente em um quarto na casa do
  \emph{melamed} (professor) e as aulas reuniam crianças das mais
  variadas idades.}--- e ali iniciam um debate as mais diversas gerações
de estudiosos e intérpretes dessa tradição. Moisés é lido ao lado de
Rabi Akiva, assim como Baal Schem Tov ao lado de Rabi Iehiel Mihal --- a
hierarquia temporal é suspensa e as noções são passíveis de
interpretação e de adaptação. E finalmente o aluno é posto ao lado do
mestre e vice-versa, até que ele próprio se torne o mestre.

É a partir do preceito de que --- na tradição judaica --- qualquer
indivíduo é convidado a fazer um \emph{chidush} que se constrói o
entendimento de que, na falta de lugares santos --- especificamente
físicos e ao mesmo tempo judaicos ---, existem milhares de tempos
santos. Como já dizia o rabino americano Abraham Hessel, ``o judaísmo constrói catedrais no tempo''. O conceito de ``lugar
santo'' --- tangível e intangível --- foge de arquiteturas estáveis e
intransportáveis.

Os lugares santos judaicos são os símbolos, os signos, os objetos
rituais, litúrgicos e religiosos que sacralizam e santificam um
território-corpo --- em suas mais variadas escalas --- e delimitam uma
silhueta de pertencimento, mesmo que temporária. Por fim, estes lugares
santos são artifícios utilizados em momentos de culto religioso e de
reza, e vinculados a uma experiência e a uma existência de tempos
religiosos, tais quais as rezas diárias e os festejos semanais e anuais.

Em \emph{Derrida, um egípcio --- o problema da pirâmide judia}, o
filósofo alemão Peter Sloterdijk, a partir da narrativa bíblica,
discorre a respeito da influência de uma política egípcia na construção
de uma identidade, de uma religiosidade e de uma espiritualidade
propriamente judaicas. Não à toa são os egípcios escolhidos como
precursores dessa suposta identidade por meio da aglomeração, da
segregação e da exclusão. Pois, para que judeus e judias pudessem ser
considerados um povo --- e isso se dá por meio do recebimento das Tábuas
da Lei ---, este conglomerado de pessoas escravizadas --- filhos dos
patriarcas Abraão, Isaac e Yaacov e das matriarcas Sarah, Rivka, Rachel
e Leah --- foi posto em liberdade por Moisés, vagueou por quarenta anos
no deserto do \emph{Sinai}\footnote{Após a saída do Egito, o recebimento
  das Tábuas da Lei e logo antes da entrada na Terra Prometida, o povo
  judeu enviou para Canaã doze espiões para investigar e colher
  informações a respeito da nova localidade em que se fixariam. Conta a
  história que dez dos doze espiões, ao voltarem, suplicaram ao povo
  hebreu para que não entrasse na terra, inundada de gigantes, feras e
  bestas. O povo, ao receber estas informações, desesperou-se e clamou
  para Deus que preferiam estar escravizados, em retorno ao Egito. Nesse
  momento, Deus entende que para entregar a liberdade a um povo
  escravizado, esses indivíduos --- a geração que foi ao mesmo tempo
  escravizada e posta em liberdade --- deveriam morrer e dar espaço ao
  nascimento de uma nova geração de pessoas já livres. Assim, é exigido
  de Moisés que conduza o povo em retorno e que vagueie com ele por
  quarenta anos no deserto do \emph{Sinai}.} e finalmente criou para si
uma divindade que ultrapassa os monumentos fixos e intransponíveis de
seus perseguidores. O judaísmo inventou para si um Deus portátil que
habita o texto e o livro.

\begin{quote}
O Egito, dessa maneira, se torna o lugar da pirâmide, evidência
incontornável da materialidade do signo, a representação do divino em
irremovíveis e inabaláveis monumentos de pedra. É o contrário do Deus
inventado pelos judeus, um Deus portátil que habita o livro, um Deus não
mais representado por monumentos intransponíveis (pirâmides) e sim por
documentos (o texto sagrado). Esta nova versão de Deus joga inesperada
luz sobre toda uma problemática ligada a meios, transportabilidades,
migrações. Se não é possível transportar deuses que moram em pirâmides
ou templos colossais, é bem mais fácil transportar um Deus que habita o
texto por mais sagrado que seja\footnote{TELLES, \emph{Posto de
  observação: reverberações psicanalíticas sobre cotidiano, arte e
  literatura}, p.\,345.}.
\end{quote}

Nesse sentido, ``meios, transportabilidades, migrações''\footnote{\emph{Ibidem}.}
adquirem uma importância religiosa até então insuspeita. Ou seja, as
diversas tentativas de construção de identidades e de comunidades
judaicas consolidam e incorporam em si as paisagens, situações e
referências daqueles que os humilham e oprimem. Em parte, é a partir de
geografias repressivas que a identidade e as comunidades judaicas
inventam, criam e constroem para si espaços sacralizados --- em suas
múltiplas escalas --- que ultrapassam seus perseguidores e, finalmente,
que se consolidam e florescem.

\begin{center}\adforn{49}\end{center}

\noindent{}A incorporação e a consolidação de traços identitários-comunitários, a
partir de geografias de pertencimento/\,exclusão, podem ser observadas em
meio ao surgimento dos primeiros guetos judaicos na Veneza
renascentista, durante o século \textsc{xvi}. Pois, como nos lembra o historiador
e sociólogo Richard Sennett:

\begin{quote}
Ao acuarem os judeus para o gueto, os venezianos acreditavam estar
isolando o mal que infectara a comunidade cristã. Eles sentiam medo de
tocar os corpos impuros\footnote{Na Veneza renascentista, a comunidade
  cristã acreditava que judeus e meretrizes eram portadores de corpos
  impuros, em que haveria uma prostituição do corpo-carne em nome da
  usura e do dinheiro. Dessa maneira, ambos os grupos sociais eram
  considerados inferiores à moral cristã, sujeitos à exclusão e
  discriminação. Essa relação também é posta em debate ao se pensar a situação
  geográfica de judeus e prostitutas no caso da zona do meretrício no
  bairro do Bom Retiro, na cidade de São Paulo no início do século \textsc{xx}. A
  delimitação de uma determinada zona da cidade em que haveria o
  trânsito de judeus recém-imigrados e prostitutas judias --- as polacas
  --- tensiona a criação de uma narrativa moral para e pela própria
  comunidade judaica. Essa narrativa --- em constante disputa pelas
  comunidades --- exclui a presença das prostitutas judias do bairro e
  da história comunitária, e até da cidade de São Paulo.} que
identificavam como vícios corruptores\footnote{\textsc{sennett}, \emph{Carne e
  Pedra}, p.\,222.}.
\end{quote}

A montagem do gueto se deu no momento em que a cidade de Veneza passava
por um período de recessão econômica. Em consequência, os cristãos
acreditavam que a segregação da diferença --- judeus e demais corpos
indesejáveis, porém necessários --- devolveria à cidade a paz divina e a
dignidade\footnote{As crenças antissemitas que circulavam na época
  sustentavam que os corpos judaicos eram portadores de doenças e de
  vícios --- culpados que eram pela morte de Jesus Cristo e frutos do
  pecado e de impurezas. A comunidade cristã acreditava que a presença
  de judeus --- dos quais muitos praticariam a agiotagem --- teria
  desencadeado uma fúria divina sobre a cidade de Veneza, causando,
  desta forma, sua recessão econômica.}. Dessa maneira, pessoas judias
acabaram sendo restritas a uma geografia repressiva --- o gueto judeu
cercado, onde a vida judaica acontecia. Apesar da construção deste
território plástico\footnote{O conceito de território plástico foi
  desenvolvido pelo pesquisador israelense Eyal Weizman no livro
  \emph{Hollow Land}. Refere-se ao artifício de criação de
  territorialidades artificiais e impositivas perante uma geografia já
  pré-estabelecida. O termo designa as fronteiras fluidas e de
  características plásticas, ou seja, que podem se alargar ou diminuir,
  inflar ou desinchar conforme as necessidades, o contexto e a situação
  em que se encontram.} --- uma territorialidade segregada e excludente
---, a população judaica de Veneza ainda foi capaz de ``construir
novas formas de vida comunitária que extrapolaram as políticas de
isolamento e de conquistar algum nível de
autodeterminação''\footnote{\textsc{sennett}, \emph{op.\,cit.}, p.\,223.}.

Ao final, a geografia de Veneza protegeu os judeus. A construção desta
territorialidade pelos canais fez com que o gueto ficasse isolado,
permitindo que os judeus --- e a vida e a identidade judaicas --- fossem
livres para realizar seus ritos, cultos, tradições e festividades; para
construir para si uma comunidade efervescente por meio da consolidação
de uma identidade compartilhada e de espaços de convívio comunitário
(como os centros comunitários e as sinagogas). Foi a partir de uma
geografia repressiva que os judeus transformaram esta
circunstância-território em uma possibilidade.

Dessa maneira, entende-se que a vivacidade das comunidades judaicas
dentro dos guetos --- seja o gueto de Veneza, o gueto de Varsóvia ou de
diversas outras geografias excludentes ---, durante os momentos de
celebração, de perseguição e de lamento, nos espaços de trânsito, se dá
a partir de constantes movimentos de transformação do espaço maldito ---
o espaço da perseguição --- em um lugar sagrado (\emph{kadosh}),
passível de apropriação, de cultivo das fagulhas pousadas e de
identidades.

Na língua hebraica, a palavra \emph{kadosh} possui dois significados. O
primeiro expressa literalmente ``separar'' ou ``apartado'', o segundo
conota sacralidade, o ``espaço e/\,ou objeto sagrado''. \emph{Kadosh},
logo, simboliza a prática judaica de resguardar momentos e objetos que
são preservados para situações extraordinárias ou de grande valor
simbólico e religioso. A ação de transformar ``\emph{arur}'' (maldito)
em ``\emph{kadosh}'' (sagrado) diz respeito a uma apropriação das
condições e situações em que alguém se vê e, assim, à criação de outros
lugares. Trata-se de decidir por si os próprios passos. \emph{Ansina,
estonses, entiendese ke, endesparte de deshar la luz pasar, es menester
de emitir luz propria}\footnote{Tradução do ladino: ``Isso, pois,
  entende-se que, para além de deixar a luz passar, é necessário emitir
  luz própria''.}.

Por fim, estes movimentos de transitoriedades, transportabilidades e
locomoções --- as constantes apropriações de territórios (geográficos ou
não) --- são observados nos monumentos judaicos e, concomitantemente,
moldam para si os conceitos de identidade e de comunidade --- bem como a
manutenção geracional destes vínculos.

\begin{center}\adforn{49}\end{center}

\noindent{}Contam os \emph{tzadikim} que uma sinagoga nunca deveria ser destruída.
Seu desmonte apenas é possível caso a comunidade que ali frequenta se
locomova --- construindo para si outro espaço comunitário, em outras
localidades. Este preceito ``joga inesperada luz''\footnote{Tomo de
  empréstimo as palavras de Sérgio Telles, \emph{op.\,cit.}, p.\,345.} no
entendimento de que, para além de arquiteturas fixas, os espaços-tempos
judaicos são construídos e consolidados por vínculos de pertencimento,
tradições, línguas, ritos e cultos que resultam em um aglomerado de
indivíduos que compartilham signos e símbolos semelhantes e se agregam,
se unem, constroem e inventam para si uma comunidade.

Ainda nesta chave, não há restrições religiosas para a construção de
espaços sinagogais --- estes locais podem ser construídos e consolidados
em qualquer localidade, contanto que possuam em seu interior as peças
primordiais para o rito religioso. Dentre elas estão: \emph{Aron
HaKodesh}\footnote{\emph{Aron HaKodesh} (em português: ``Arca Sagrada'')
  é uma espécie de armário em que se guarda os textos bíblicos e os
  pergaminhos sagrados, a Torá. Este é um objeto posicionado
  sempre em direção ao Templo de Jerusalém.}, \emph{Bimá}\footnote{\emph{Bimá}
  (em português: ``Altar'') é uma mesa ao centro da sinagoga, usada para
  o suporte de leitura de livros e de textos bíblicos, além de ser o
  local onde o \emph{Chazan} (condutor da oração) se posiciona.} e
\emph{Esh}\footnote{\emph{Esh} (em português: ``Fogo celestial'') é uma
  vela, um candelabro ou uma luz que representa a presença divina na
  casa de orações. Assim como era feito no Templo Sagrado e lembrado durante a
  festividade de \emph{Chanucá}.}. Estes objetos, de cunho religioso,
para além de delimitar uma espacialidade judaica, são móveis,
transitórios e, dessa forma, ocupam, consagram e sacralizam lugares ---
em suas mais variadas escalas ---, animam estas territorialidades,
instalam nelas uma alma.

Por fim, os lugares santos judaicos, aqui investigados, se manifestam na
escala do corpo (\emph{kipá}), da casa (\emph{mezuzá}), da cidade
(\emph{eruv}) e do cosmos (\emph{makom}); cosmogonizam um território
(real, ficcional e/\,ou imaginário) e, para além de lembretes diários,
firmam-se enquanto declarações públicas, individuais, comunitárias de
identidade, de crença e de fé.

\begin{center}\adforn{49}\end{center}

\noindent{}A \emph{kipá}\footnote{A \emph{kipá} é um chapéu, boina ou touca usado
  tradicionalmente por homens e recentemente incorporado por mulheres
  nos movimentos religiosos judaicos reformistas e reconstrutivistas. É
  um objeto que delimita um território-corpo de pertencimento e se firma
  enquanto um lembrete constante da presença de Deus.} e a
\emph{mezuzá}\footnote{A \emph{mezuzá} é um cilindro de metal, de
  madeira ou de acrílico posto nos batentes das portas judaicas e que
  contém um pergaminho com a bênção do \emph{Shemá Israel}. Este ato é
  uma tradição e um lembrete que remete aos tempos de escravização no
  Egito, tendo sido reinterpretado pelas comunidades. Judeus
  sefaraditas costumam posicionar a \emph{mezuzá} de forma
  paralela ao batente, enquanto judeus ashkenazitas a posicionam
  de forma inclinada.} são eixos cósmicos, signos de ligação e de
lembrança da presença divina e de sua conexão com o mundo humano. Estes
símbolos são objetos hierofânicos --- manifestações da realidade sagrada
de Deus por meio de artifícios e de signos do mundo real, e assim
passíveis de compreensão pelas interpretações humanas. Eles se
consolidam enquanto atos misteriosos e distinguem a realidade sagrada de
um mundo profano. Pois ``para o homem religioso o espaço não é
homogêneo\footnote{Em \emph{A poética do espaço}, o filósofo e poeta
  francês Gaston Bachelard aponta que ``nem o tempo, nem o espaço
  em que vivemos são vazios, neutros e homogêneos, mas são carregados de
  qualidade e habitados por fantasmas: os espaços de nossos devaneios e
  paixões podem ser leves, etéreos, transparentes, obscuros ou
  pedregosos, móveis e fixos'', enfim, vivemos em um conjunto de
  relações que, segundo Bachelard, definem nossas experiências,
  agregações individuais e coletivas, e suas possibilidades de
  reinvenção. \textsc{bachelard} \emph{apud} \textsc{rago}, \emph{Inventar outros espaços, criar
  subjetividades libertárias}, p.\,14.}, o espaço apresenta
roturas, quebras; há porções de espaço qualitativamente diferentes de
outros''\footnote{\textsc{eliade}, \emph{O sagrado e o profano}, p.\,25.}.

Neste sentido, a presença divina transcende o mundo espiritual --- em
direção ao mundo material/\,humano --- por meio de objetos humanos
conhecidos e que fazem parte de suas próprias realidades. Por exemplo,
no momento em que a divindade revelou-se a Moisés a partir de uma árvore
em chamas, o objeto ``árvore'' não é adorado como árvore em si, mas sim
enquanto uma hierofania --- um dispositivo que permite a passagem do
sagrado ao mundo humano-profano, já que, ao final, o sagrado se
manifesta quando revela-se.

\begin{quote}
Não te aproximes daqui, disse o Senhor a Moisés; tira as sandálias de
teus pés. Porque o lugar onde te encontras é uma terra Santa\footnote{ÊXODOS,
  3:5.}.
\end{quote}

Assim sendo, ``toda hierofania espacial ou toda consagração de um
espaço equivale a uma cosmogonia''\footnote{\textsc{eliade}, \emph{op.\,cit.}, p.\,59.}.
A \emph{kipá} --- este eixo cósmico de ligação entre o ser humano e o
divino ---, neste sentido, revela-se enquanto consagração de um
determinado território-corpo, um objeto hierofânico portátil, usado por
judeus e por judias como revelação e recordação da presença divina no
mundo --- em que, para além de suas próprias silhuetas, está firmada a
divindade, criadora do mundo e de todas as coisas que ele contém. O eixo
cósmico --- a ligação direta com Deus --- é reforçado a partir da
utilização deste artifício-objeto, pois, ao final, segundo o escritor e
teólogo Mircea Eliade\footnote{\emph{Ibidem}, p.\,35.}, o ``homem
religioso'' deseja estar a todo instante em ``Seu Mundo'', em contato
direto com o divino.

Para além de um objeto hierofânico, a \emph{kipá} se apresenta enquanto
um marcador legal, social, identitário e comunitário de pertencimento
judaico. A utilização deste artifício-objeto se estabelece enquanto um
delimitador de pertencimento e de congregação entre os diversos
indivíduos que transitam em um determinado território onde a comunidade
se estabelece; tais indivíduos, por sua vez, se reconhecem como
portadores de uma identidade compartilhada. E assim, essa peça delimita,
demarca e consagra um corpo/\,\emph{corpus}, possibilitando um senso de agregação
entre os transeuntes de uma determinada territorialidade-cidade onde a
comunidade decide fixar-se.

O uso da \emph{kipá} firma-se enquanto uma possibilidade de
``reconhecer-se judeu ou judia'' em comunidade --- para além do próprio
corpo e da conexão religiosa com o divino. Permanecer em comunidade
permite ao ``homem religioso'' --- este \emph{corpus} de pessoas ---
estar a todo instante em ``Seu Mundo'' --- ou seja, em contato direto
com Deus e com seus iguais. Dessa maneira, o ``homem religioso'', ao
utilizar o artifício dos objetos hierofânicos, cosmogoniza um território
comunitário e cria, por fim, um corpo, um local e um mundo santificado
para si.

A cosmogonização e a consagração de um mundo --- o ``Seu Mundo'' --- são
compreendidas a partir do conceito de eixo cósmico --- uma relação
direta e vertical entre o ser humano e o divino. Segundo Eliade, é a
partir do eixo cósmico --- e em volta dele --- que o território se torna
habitável pela presença divina. Ou seja, é a partir da consagração
destas territorialidades que o ``homem religioso'' faz do espaço profano
--- amorfo e sem limites --- o ``Seu Mundo'', pois ``Nosso Mundo é
a `terra santa' porque é o lugar mais próximo de Deus''\footnote{\emph{Ibidem},
  p.\,40.}.

``A instalação de um território equivale à fundação de um
Mundo''\footnote{\emph{Ibidem}, p.\,46.}. Essa instalação --- ligada ao eixo
cósmico --- ritualiza, sacraliza e demarca um espaço, em suas mais
variadas composições. Parte-se do pressuposto de que o ``homem
religioso'' necessita estar a todo instante em contato com Deus; assim
sendo --- para além de seus corpos ---, estes indivíduos desejam que
suas próprias casas e moradas se tornem o ``Centro do Mundo''. A morada
--- a habitação do ser humano --- situa-se no centro de seu próprio
mundo religioso e reproduz, em escala microcósmica, o Universo criado
por Deus.

Ao utilizar o artifício-objeto hierofânico da \emph{mezuzá}, famílias
judias (e as comunidades) reafirmaram a presença de Deus em suas
próprias moradas. A casa e/\,ou o espaço comunitário (sinagoga), por fim,
se tornam uma \emph{imago mundi}\footnote{\emph{Imago mundi} é um termo
  do latim que representa a ``imagem/\,espelhamento do mundo criado por
  Deus''.} --- ou seja, ao passar por um ritual de consagração e
sacralização, tornam-se um espaço religiosamente habitável, em conexão
com Deus a partir do eixo cósmico. Nos textos talmúdicos, a casa
judaica religiosa é referida como \emph{Bait Neeman} --- uma ``casa
leal'' --- uma habitação santificada pela presença divina, já que ela
própria reitera a criação do mundo e a fidelidade e dedicação ao
monoteísmo e ao Deus judaico.

Essa tradição remete ao momento em que Moisés suplica ao Faraó a
liberdade de seu povo escravizado no Egito. Segundo a narrativa bíblica,
a restrição dessa liberdade impõe que Deus envie ao Egito dez pragas,
sendo a última delas a descida à terra do Anjo da Morte e o assassinato
de primogênitos. Para que os hebreus fossem poupados, Moisés aconselha o
povo para que pinte os batentes das portas com sangue de carneiro. Esse
ritual consagra e delimita um território de crença e de devoção a Deus
--- santifica uma habitação, a torna ``fiel'' à presença divina.

Enquanto uma reinterpretação do sangue de carneiro despejado pelos
hebreus nos umbrais de suas portas, a \emph{mezuzá} é um artifício de
demarcação espacial que eleva um território-casa. É uma caixa de
acrílico, metal ou madeira fixada nos batentes das portas judaicas e que
contém um pergaminho com a bênção do \emph{Shemá Israel}\footnote{A reza
  de \emph{Shemá Israel} é uma das mais importantes bênçãos judaicas;
  afirma a lealdade a Deus e ao monoteísmo.}. Este objeto é portátil,
removível, ajustado e acompanha as comunidades nos momentos de diásporas
e dispersões, voluntárias ou forçadas. A \emph{mezuzá} reitera a
cosmogonia e a devoção judaica de lealdade a Deus, santifica uma
habitação e se estabelece enquanto um simbolismo cósmico, um reflexo da
própria criação divina do mundo. Cria-se, a partir dela, o eixo cósmico.
Já que esse objeto é a chave da casa, fixar-se em um local é instalar
nele, quase que instantaneamente, uma \emph{mezuzá}.

Se toda criação humana repete a fundação do mundo, para que esta criação
dure --- perdure no espaço e no tempo --- ela necessita ser animada,
receber vida, ``uma alma''\footnote{\textsc{eliade}, \emph{op.\,cit.}, p.\,53.}. Segundo
Eliade\footnote{\emph{Ibidem}, p.\,11.}, ``para o homem primitivo tudo é
dotado de alma''. Em uma sinagoga, em um centro religioso ou
comunitário, a comunidade que ali transita é a alma e a chama de vida
dessa localidade (o \emph{minian}). A comunidade --- essa congregação de
indivíduos em prol de objetivos em comum --- funda um território de
pertencimento que anima o lugar --- introduz nele uma alma. Assim sendo,
as comunidades judaicas transformam os espaços profanos em lugares
sagrados por meio da consagração destes territórios. Nestas
territorialidades tudo é dotado de vida, de celebração e de rito, tudo
parte de uma orientação. O ``homem religioso'', por fim, transforma o
maldito em sagrado --- constrói suas identidades, santifica esses
espaços e os compartilha.

Instalar-se em um determinado lugar é reiterar a cosmogonia --- imitando
a obra divina de criação do mundo por Deus. ``Fixar-se em uma
região é assumir responsabilidade por este mundo criado; mantê-lo vivo e
renová-lo''\footnote{\emph{Ibidem}, p.\,54.}. Fixar-se em um determinado local
equivale a um novo começo. E, assim como um calendário cíclico, todo
começo repete o início primordial da criação divina do mundo.

Já que instalar-se em uma determinada territorialidade funda de novo um
mundo, essas tentativas de cosmogonização são formas de dedicar-se a
viver o cosmos puro e santo, tal como era no princípio, quando saiu das
mãos da divindade. Estas propostas são observadas nos objetos
hierofânicos --- \emph{kipá} e \emph{mezuzá} --- e relacionam-se à
experiência de tempos santos judaicos, cíclicos, elipsoidais e
constantes. A vitória do lugar sagrado sobre o espaço profano deve ser
repetida simbolicamente a todo instante, pois a todo instante o mundo
necessita ser criado novamente desde o momento primordial.

Ao fim e ao cabo, a experiência do sagrado funda ontologicamente o
mundo. ``O espaço sagrado tem um valor existencial para o homem
religioso. Nada se pode começar, nada se pode fazer sem uma orientação
prévia''\footnote{\emph{Ibidem}, p.\,26.}. Este ponto fixo --- a criação do
mundo --- permite uma direção. Se, para viver no mundo, o ``homem
religioso'' necessita fundá-lo, nenhum mundo pode nascer do ``caos'', da
homogeneidade e da relatividade do espaço profano.

\begin{center}\adforn{49}\end{center}

\noindent{}Fundar um mundo, orientar-se em meio ao ``caos'' --- ao relativismo do
espaço profano ---, torná-lo sagrado, é sempre possível. Para isso pegue
o primeiro texto que encontrar, \emph{texto-memória},
\emph{texto-qualquer}, um \emph{texto-nenhum}. Ele sempre será um mapa.

Estes monumentos móveis são Deus, que habita o texto e o livro. São
faróis e dispositivos de direção que acompanham as comunidades nos
constantes movimentos das diásporas --- pousam e fixam-se em outras
territorialidades; inventam e fundam outros, mas nunca novos, mundos.

\chapter{Geografias}

\begin{quote}
``Um hassid perguntou ao Maguid de Zlotschov. Nosso mestre Rashi diz: \emph{O
que faltava ao mundo adormecido? Só o descanso. Veio o \emph{shabat} e, com
ele, o descanso}. Por que não se diz simplesmente que ao mundo faltava o
descanso, até que veio o \emph{shabat}? Pois se \emph{shabat} e descanso querem dizer
a mesma coisa!

--- \emph{Shabat} --- disse o Rabi --- quer dizer a volta ao lar, porque neste
dia todas as esferas voltam ao seu verdadeiro lugar. É para isto que
aponta Rashi: \emph{inquietas estão as esferas durante a semana, porque foram
descidas de seus lugares, e no \emph{shabat} encontram o descanso, porque podem
voltar}''\footnote{\textsc{buber}, \emph{op.\,cit.}, p.\,196.}.
\end{quote}

\noindent{}Segundo a filosofia judaica, o \emph{shabat} é o sétimo dia da criação
divina, momento em que a divindade apenas contemplou o mundo e suas
coisas. Como recordação da criação do mundo por Deus, o \emph{shabat} é
considerado um período sagrado e dedicado exclusivamente ao repouso.
Inicia-se no pôr do sol da sexta-feira e termina ao nascer das três
primeiras estrelas do sábado, totalizando 25 horas de duração.

Durante o \emph{shabat}, judeus e judias recordam o ato de criação
divina do mundo --- e de todas as coisas que ele contém --- por meio do
descanso físico, mental e espiritual e da pausa do trabalho criativo. O
momento do \emph{shabat}, dessa forma, se mostra enquanto um espaço (e
um tempo) em que se permite o ócio, propício a reflexões, ``uma
volta ao lar'', segundo o Maguid de Zlotschov --- e assim a conexão com
a divindade e com a comunidade\footnote{As comunidades judaicas se unem
  e se apropriam de um território real para criarem seu mundo e
  realizarem seus ritos religiosos.}. Afinal, ``sem o ócio não
poderia haver criação, apenas repetição monótona do mesmo, submissão à
esfera da necessidade''\footnote{\textsc{rago}, \emph{op.\,cit.}, p.\,10.}.

Nessa lógica, são ao todo 39 \emph{melachot}\footnote{\emph{Melachot}
  (plural de \emph{melachá}) são ações criativas que interferem
  diretamente na lei do descanso judaico durante o \emph{shabat}. As
  \emph{melachot} são atividades que alteram a ordem natural das coisas,
  bem como a maneira como se encontrava o mundo no exato momento em que
  foi concebido o \emph{shabat}.} proibidas durante o rito do
\emph{shabat}, sendo elas: plantar, arar, cortar, juntar, debulhar,
dispersar, separar, moer, peneirar, fazer massa, assar, tosquiar, lavar,
desembaraçar, tingir, fiar, esticar, passar, tecer, desfazer nós e/\,ou
fios, desatar, costurar, rasgar, caçar, abater, pelar o couro, curtir o
couro, demarcar o couro, cortar, escrever, apagar, construir, destruir,
apagar o fogo, acender o fogo, terminar a manufatura de um objeto,
transportar de uma propriedade à outra ou carregar em um domínio
público.

Durante o \emph{shabat}, carregar objetos ou corpos, entre domínios
públicos\footnote{Segundo as interpretações rabínicas, domínios públicos
  ou, em hebraico, \emph{reshut harabin}, são todos aqueles espaços-ruas
  que conectam uma cidade de uma ponta a outra, que tenham pelo menos
  doze metros de largura e onde transitem 600 mil pessoas por dia. Um
  exemplo de domínio público na cidade de São Paulo seria a avenida
  Paulista.} e domínios privados\footnote{Segundo as interpretações
  rabínicas, domínios privados ou, em hebraico, \emph{reshut hayachid},
  são todos aqueles espaços que tenham pelo menos 80 cm de elevação do
  solo e uma área mínima de 4 por 4 \emph{tefachim}, aproximadamente 40
  x 40 cm, podendo um indivíduo e/\,ou associação ser proprietária desse
  local, ou não.} (e vice-versa), é proibido segundo a
\emph{Halachá}\footnote{\emph{Halachá} é um conjunto de normas e leis
  que servem como guia para uma vida judaica religiosa. A palavra
  \emph{Halachá} vem do vocabulário hebraico, mais especificamente do
  termo \emph{heh-lamed-kaf}, que significa viagem e/\,ou caminho.}.
Assim, é dentro dessa perspectiva que surge o \emph{eruv}. Um artifício
que se estabelece enquanto um portal nas territorialidades em que as
comunidades se agregam, ou seja, ``uma linha que permite a junção
de vários tempos e espaços, e que ali, neste ponto para o qual tudo
converge, seria possível o trânsito entre passado e futuro, sonho e
vigília, vida e morte''\footnote{\textsc{saavedra}, \emph{O manto da noite}, p.\,86.}.

O \emph{eruv}, por assim dizer, é um enclausuramento ritualístico que
permite a abertura dos espaços privados da morada ao domínio público ---
criando, desta forma, uma casa aberta com becos, vielas e pátios
compartilhados entre a comunidade --- permitindo a circulação e o
transporte de objetos.

A fantasia do \emph{eruv} acontece às escondidas. Este território mágico
surge nas frestas e nas brechas de uma proibição judaica religiosa --- o
carregamento de objetos durante o rito judaico do \emph{shabat}. O
\emph{eruv} materializa o texto --- e suas interpretações --- no espaço
e se justifica a partir dos debates rabínicos sobre os livros sagrados.
O \emph{eruv} é uma estratégia que contesta as restrições\footnote{Mas a
  partir das regras do jogo --- ou seja, a partir das interpretações
  rabínicas dos textos.} de carregamento, e mais que nada entende que,
para além das práticas religiosas, as comunidades judaicas sobrevivem em
diásporas a partir de suas agregações, seus \emph{minianim}\footnote{\emph{Minianim}
  (plural de \emph{minian}) refere-se ao quórum mínimo de dez homens
  judeus adultos necessários para realizar certas obrigações religiosas,
  como a retirada da Torá da Arca Sagrada e sua leitura. Mais
  recentemente e em algumas vertentes religiosas e liberais judaicas, as
  mulheres também passaram a contabilizar como participantes ativas do
  \emph{minian}.}, seus laços e traços de pertencimento e de
sociabilidade.

Para a pesquisadora e arquiteta Miriam Levy, ``o \emph{eruv} é uma forma
de navegação nômade sobre a topografia de um espaço, em que a noção de
pertencimento a uma comunidade é criada e construída em detrimento de um
território em particular''\footnote{\textsc{levy}, \emph{Encounter with the
  \emph{eruv}: a project toward the city of open enclosures}, p.\,51. (tradução
  nossa)}. Por fim, independentemente deste local, o \emph{eruv} pode
ser acionado e construído em qualquer lugar e em qualquer tempo, desde
que haja uma comunidade agregada.

\begin{center}\adforn{49}\end{center}

\noindent{}\emph{Eruv} provém da palavra aramaica \emph{arab}, que em português
significa ``mistura''. Ou seja, o \emph{eruv} é um artifício judaico que
permite a combinação de espaços públicos e espaços privados,
apropriando-se destas territorialidades para criar assim um outro local
de pertencimento, de trânsito e de circulação. De forma prática, este
artifício se estabelece no tecido urbano das cidades por meio de um
\emph{tsurat hapetach} (em português: ``formas de portas''), dois postes
verticais metálicos implementados paralelamente às calçadas, rentes às
construções da cidade, e conectados em suas extremidades por um fio de
nylon transparente ou por um barbante.

A marcação deste espaço físico surge enquanto resposta às restrições de
carregamento no \emph{shabat}. Assim, o \emph{eruv} é nômade e
construído por meio de conexões e interferências com os espaços urbanos
e suas edificações pré-existentes. Em um contexto diaspórico há uma
necessidade de criar e consolidar uma casa comum e compartilhada em um
determinado território --- apropriar-se dele, consagrá-lo desde dentro.
Por meio desta estrutura física e simbólica permite-se a solidificação
de espaços judaicos em que a comunidade se conecta com uma cultura
pré-existente e, simultaneamente, também com a sua própria. E assim,
``o \emph{eruv} faz parte da topografia urbana da cidade e interage com
ela''\footnote{\emph{Ibidem}, p.\,13. Tradução da autora.}.

O \emph{eruv} possibilita o estreitamento de laços dentro da comunidade
que está inserida em um determinado território por ele delimitado. A
apropriação destes espaços molda as dinâmicas sociais, religiosas,
simbólicas e seus deslocamentos durante o rito do \emph{shabat}.
Entende-se aqui que o \emph{eruv} é um artifício de pertencimento por
meio do qual indivíduos judeus se locomovem nos territórios em que se
fixam.

Estes deslocamentos entre espaços públicos e espaços privados permitem
uma identificação e uma agregação entre os membros da comunidade
judaica; ao final, o contato visual nas ruas entre minorias e com uma
cultura maior (sociedade) estreita o senso de comunidade, de
identificação, e ajuda os indivíduos a formar laços políticos e alianças
sociais de resistência nos territórios diaspóricos.

Dessa forma, ao sair às ruas, a comunidade torna o ambiente externo
menos ameaçador, faz dele uma localidade apropriada. Ao final, o senso
de pertencimento e as agregações comunitárias são atos de resistência
que demarcam e delimitam territórios, solidificam as relações
interpessoais e apresentam a comunidade enquanto um coletivo.

\begin{center}\adforn{49}\end{center}

\noindent{}O \emph{eruv} é uma área circunscrita por meio da união de construções
pré-existentes e de fios --- paredes físicas e paredes simbólicas,
respectivamente --- em que se permite o enclausuramento e a apropriação
destas territorialidades. É por meio da apropriação deste local que se
abre uma ruptura no espaço-tempo, possibilitando relações de
sociabilidade durante o rito do \emph{shabat}. Assim, como forma de
cumprir as restrições de carregamento do \emph{shabat}, os espaços
urbanos, supostamente laicizados, são alterados e ressignificados pelas
ações e pelas interações e interpretações religiosas.

Neste sentido, o \emph{eruv} se torna uma casa-território compartilhada
pela comunidade judaica local durante o rito do \emph{shabat}. É uma
territorialidade em que os espaços coexistem, porém não mais coordenados
pelo binômio público-privado\footnote{O Talmud divide os espaços
  e as territorialidades em duas categorias: \emph{reshut harabin} e
  \emph{reshut hayachid} (espaços públicos e espaços privados,
  respectivamente). Os rabinos, a partir das interpretações das
  características espaciais, definem mais dois lugares, \emph{carmelit}
  e espaços neutros. Os espaços \emph{carmelit} são locais que
  apresentam características de espaços públicos, porém sem as dimensões
  exatas necessárias, por exemplo, rua de bairro, vila de casas. Já os
  espaços neutros são territórios que se assemelham aos espaços
  privados, mas que possuem menos de 1,6m² de área ou 80 cm de altura,
  por exemplo, postes de eletricidade, decks de piscinas, tábua de uma
  mesa.}. Segundo Miriam Levy, ``é nestes locais que é possível
encontrar outros sinais de identidade e novos espaços de
colaboração''\footnote{\textsc{levy}, \emph{op.\,cit.}, p.\,21. Tradução da autora.}.

O \emph{eruv} é uma mescla de pátios. Este encontro de domínios, espaços
públicos e espaços privados apropriados se dá por meio de uma refeição
compartilhada entre a comunidade. Pois, no caso judaico, a comunidade é
um marcador legal e social de pertencimento. Todas as pessoas que ali se
alimentam fazem parte de uma mesma comunidade\footnote{Os primeiros
  \emph{eruvim} (plural de \emph{eruv}) surgem nas comunidades judaicas
  logo após a destruição do segundo Templo Sagrado de Jerusalém (76
  d.E.C.), como uma tentativa de criação de comunidade para além das
  restrições do \emph{shabat}. Este convívio comunitário era realizado
  por meio de uma associação entre os diversos moradores das casas, que
  compunham uma vila, por meio de uma refeição compartilhada. No momento
  do \emph{shabat}, as vilas judaicas eram cercadas por tábuas de
  madeira posicionadas nos topos das construções, e assim simbolicamente
  criavam um umbral/\,porta de pertencimento à localidade.}. Essa
demarcação constrói uma casa simbólica no território da cidade, e é
entoada por meio da bênção do \emph{eruv}:

\begin{quote}
Este é o \emph{eruv} que deve ser utilizado para permitir que nós possamos
trazer coisas de dentro para fora. De nossas casas para os nossos
quintais. De nossos quintais para nossas casas. De casa em casa, de
quintal em quintal, de terraço em terraço. De nossas casas e nossos
quintais para becos e vielas, e de becos e vielas para suas casas. Para
nós, e para todas as pessoas que vivem aqui, e quem se juntar conosco
durante o rito do \emph{shabat} e nas festividades\footnote{Bênção judaica do
  \emph{eruv} entoada sobre um alimento que será repartido entre a
  comunidade.}.
\end{quote}

``A comida'', como esclarece Levy, ``é um símbolo de que
não importa em que localidade estas pessoas estejam; porém, entende-se
que esses indivíduos estão em um determinado espaço e tempo em que está
também a comunidade, e assim, onde a comunidade se
estabelece''\footnote{\textsc{levy}, \emph{op.\,cit.}, p.\,59. Tradução da autora.}. Na
maior parte dos casos, essa comida compartilhada em comunidade é uma
\emph{matzá}, um pão ázimo sem fermentação, como preparado originalmente
durante a fuga bíblica da escravidão no Egito. Simboliza, ao mesmo
tempo, a transitoriedade --- o pão que não ficou pronto --- e o
pertencimento.

É a \emph{matzá} --- essa refeição --- que é expandida simbolicamente no
território físico e urbano da cidade, sendo ela também o próprio
\emph{eruv} --- possibilitando que toda a comunidade compartilhe este
alimento. Os fios e postes, ao final, são delimitadores físicos (pouco
visíveis, se não invisíveis para muitos) no território urbano --- a
partir deles entende-se quais ruas estão englobadas e quais ruas estão
descartadas dentro da marcação do ``Distrito de
\emph{Eruv}''\footnote{Apesar do \emph{eruv}, em grande maioria, se
  apresentar enquanto uma construção fictícia, imaginária e silenciosa
  de bairro e de agregação comunitária, a consolidação do \emph{eruv} na
  cidade de São Paulo necessitou de um contrato legal e jurídico entre a
  comunidade judaica e a prefeitura --- esta permite que a comunidade
  alugue simbolicamente os direitos do Tesouro Municipal, para, dessa
  forma, implementar os postes nas calçadas públicas da cidade. O
  contrato, assim sendo, garante que a comunidade judaica seja
  legalmente (e simbolicamente) proprietária do território no qual o
  \emph{eruv} é instalado --- chamado de ``Distrito de \emph{Eruv}''.}.

O \emph{eruv} é um marcador legal e social de pertencimento comunitário.
Estar dentro desta territorialidade, uma construção fictícia de bairro
--- inventado e apropriado pelas comunidades judaicas ---, permite uma
simbolização. Um estabelecimento \emph{kosher}\footnote{\emph{Kosher}
  significa ``apropriado'' e/\,ou ``adequado'' e são preceitos e práticas
  que seguem as normas da religiosidade judaica. Um exemplo são os
  alimentos \emph{kosher}, preparados segundo as leis alimentares
  judaicas, sendo elas: carne de animais que sejam, ao mesmo tempo,
  ruminantes e de patas fendidas, peixes que possuem escamas e
  barbatanas, além da proibição da mistura de carne com laticínios.},
apesar de não funcionar durante o \emph{shabat}, necessita estar dentro
das delimitações do \emph{eruv}. Assim como estão as sinagogas, os
centros comunitários, as escolas e as casas judaicas.

O \emph{eruv}, neste sentido, é um assunto comunitário. Durante o
\emph{shabat}, há a obrigação religiosa da ida à sinagoga. Este espaço,
mais do que uma casa de orações, é um centro comunitário, de convívio,
de estudos --- e assim, de socialização. ``O \emph{eruv} é sobre a
movimentação de pessoas que estão em suas casas para as
ruas''\footnote{\textsc{levy}, \emph{op.\,cit.}, p.\,33. Tradução da autora.}. Se, por conta
das leis restritivas do \emph{shabat}, pessoas com necessidades
especiais não pudessem sair de casa, a comunidade estaria inviabilizando
uma vida judaica coletiva. Sem o \emph{eruv}, indivíduos que utilizam
bengalas, cadeiras de rodas, pessoas com mobilidade reduzida, portadores
de deficiências físicas, adultos com crianças de colo que necessitam do
auxílio de carrinhos de bebês não se deslocariam de suas casas para a
confraternização com a comunidade e, assim, ficariam privados de viver
um judaísmo em coletivo.

Nesta perspectiva, o artifício do \emph{eruv} se sustenta enquanto um
fortificador diaspórico de comunidade, no sentido mais amplo desta
palavra. Afinal, viver em comunidade, segundo o escritor e pesquisador
Luiz Antônio Simas, ``representa a agregação de diversos
indivíduos em busca de objetivos comuns''\footnote{\textsc{simas}; \textsc{rufino}, \emph{op.\,cit.}, p.\,4.}, ou seja, criar e manter conexões --- mais que nada,
entende-se aqui que pessoas com dores, e com sonhos, criam comunidades.

\begin{center}\adforn{49}\end{center}

\noindent{}``O \emph{eruv} é a construção de um único espaço privado/\,domesticado ao
redor de um bairro/\,comunidade, em que carregar objetos torna-se
permitido''\footnote{\textsc{levy}, \emph{op.\,cit.}, p.\,83. Tradução da autora.}. Assim, a
partir de leis religiosas, a estratégia do \emph{eruv} cria uma série de
ressignificações nos espaços físicos e urbanos das cidades
e ``apaga as divisões legais entre os espaços, cria um grande
espaço privado, domesticado''\footnote{\emph{Ibidem}, p.\,36. Tradução da autora.}.

As leis judaicas religiosas do \emph{shabat} permitem o transporte de
objetos e utensílios nos espaços privados, mas não entre um espaço
privado e outro público. Tal interdição justifica-se porque os objetos
--- bolsas, celulares, carteiras, livros etc. --- podem ser obstáculos a
relações humanas mais aprofundadas. Como nos lembra o Maguid de
Zlotschov, ``o \emph{shabat} é uma volta ao lar''\footnote{\textsc{buber}, op.
  cit., p.\,196.} em que as relações interpessoais são intensificadas.

O \emph{shabat} é um convite para olhar, para refletir e para se
encontrar com o outro dentro de si mesmo. Tudo no \emph{shabat} está
voltado e vinculado ao descanso mental e espiritual. Nesse período,
todos os atos, ritos e práticas judaicas se relacionam ao conceito de
\emph{shalom} (paz).\footnote{Em ``O Santo \emph{Shabat}'': ``A
  mãe do Rabi Mendel perguntou-lhe certa vez: \emph{O que significa realmente
  dizermos: o santo \emph{shabat}?}. Ele respondeu: \emph{O \emph{shabat} santifica}. Ela,
  porém, disse: \emph{Não só santifica, como também cura}'. \textsc{buber}, \emph{op.\,cit.}, p.\,446.} Este encontro com a paz parte do
convívio comunitário, no entrelaçamento entre a comunidade, o descanso,
a religiosidade e a divindade. Afinal, os seres humanos foram criados
para terem prazer com a conexão divina.

A crença central do \emph{shabat} procura desenvolver um espírito de
paz. O \emph{shabat} começa com a imagem de uma noiva e de um noivo
caminhando em direção à \emph{chupá}\footnote{\emph{Chupá} é uma tenda
  sob a qual se realiza o casamento judaico.} e termina com a presença
do profeta Eliahu, anunciador da era messiânica e da vinda do terceiro
Templo Sagrado de Jerusalém. A união dos noivos é expandida de forma a
incluir a esperança da redenção do mundo, por completo\footnote{Ao
  final, ``Aquele que estabelece a paz nas alturas, há de
  estabelecer a paz entre nós''. Jó 25:2}. O \emph{shabat} começa e termina com acendimento de
fogo. Em um primeiro momento, ao pôr do sol da sexta-feira, utiliza-se
duas velas independentes --- dois castiçais acendidos separadamente. E,
ao final do período, ao nascer das três primeiras estrelas ao sábado,
utiliza-se uma única vela de \emph{havdalá},\footnote{\emph{Havdalá} é
  uma cerimônia religiosa judaica que marca o fim do período do
  \emph{shabat}, dando início a uma nova semana.} em formato de trança
que se enrosca.

\begin{quote}
O \emph{shabat} começa com a oração que descreve o primeiro \emph{shabat} no Éden,
onde existiam apenas Adão e Eva, e termina com a esperança e a oração de
que a unicidade de Deus se reflita na unicidade do povo.\footnote{WEISS,
  \emph{A Microcosm of the Shabbat Spirit}, p.\,5. Tradução da autora.}
\end{quote}

Se a unicidade de Deus está refletida na unicidade do povo, então a
filosofia judaica parte do pressuposto de que já há uma familiaridade e
unicidade (um sentido de paz) entre os diversos indivíduos que
compartilham um determinado espaço de convivência; e assim, faz-se
permitido o transporte de objetos dentro dos espaços privados e/\,ou
domesticados --- a casa. Ao final, a casa, a morada, por si só
estabelece uma pequena comunidade, desde que haja uma certa intimidade
entre os indivíduos que nela habitam.

Assim, ao cercar uma determinada territorialidade da cidade, o
\emph{eruv} não incide diretamente em espaços públicos como os
conhecemos. O \emph{eruv} cerca e delimita espaços urbanos considerados,
segundo as interpretações rabínicas, espaços \emph{carmelit}.
\emph{Carmelit} são as ruas de bairro, pequenas vilas de casas, ruas com
comércios menos movimentados. O \emph{eruv}, ao circunscrever-se nos
espaços \emph{carmelit} --- uma territorialidade de pouco trânsito ---,
pressupõe que os indivíduos que transitam por estas zonas já fazem parte
de um microcosmo --- uma ideia de casa, de intimidade ---, a comunidade
em si já existe\footnote{Este é o caso da Festa do \emph{Eruv},
  realizada pela Casa do Povo anualmente e vista pela comunidade do Bom
  Retiro enquanto uma festa de bairro, apropriada pelos indivíduos que
  transitam, moram ou moram ali. Por assim dizer, mesmo que o conceito
  do \emph{eruv} tenha nascido dentro de um ambiente judaico, a Casa do
  Povo, durante a Festa do \emph{Eruv}, traduz este conceito e engloba
  todas as pessoas, transformando o espaço da Casa em um território
  compartilhado. Assim, ao se apropriar do espaço público da rua Três
  Rios (local onde está situada a Casa do Povo no Bom Retiro), a Festa
  do \emph{Eruv} convida toda a comunidade do bairro a frequentar este
  território. Durante o evento diversos são os indivíduos que participam
  ativamente na organização, nas oficinas e nas atividades propostas.
  Afinal, o Bom Retiro é conhecido como um bairro multicultural, em que
  transitam e onde se estabelecem diversos grupos sociais: judeus,
  coreanos, paraguaios, bolivianos, venezuelanos, nordestinos e
  comunidades negras, pessoas LGBT+, prostitutas, gregos. A casa se
  torna um espaço comum, uma outra frequência.}.

Se partirmos do pressuposto de que o \emph{eruv} combina a arquitetura
local com uma arquitetura judaica religiosa --- um espaço (e um tempo)
físico e imaginário, real e ficcional, concreto e abstrato ---, então
ele constrói para si um ``espaço de pertencimento que é nômade em
sua forma de construção, mas estacionário em suas conotações
metafóricas''\footnote{\textsc{levy}, \emph{op.\,cit.}, p.\,61. Tradução da autora.}. O
\emph{eruv}, paradoxalmente, cria um espaço em que há um sentimento
comunitário de pertencimento sedentário, até que ele se desmanche ---
``este local de pertencimento se move com a comunidade, quando
necessário''\footnote{\emph{Ibidem}, p.\,59. Tradução da autora.} --- e seja
construído em outro território.

\begin{center}\adforn{49}\end{center}

\noindent{}O \emph{eruv} é uma fronteira, um limite e um limiar em torno de toda
uma comunidade e se estabelece enquanto uma multiplicidade de espaços e
de tempos sagrados e multifacetados. Este espaço sagrado (e apropriado)
engloba pessoas, objetos, animais, textos, livros e atividades
ritualísticas, religiosas, cotidianas, coletivas e comuns, bem como as
suas arquiteturas. Durante o \emph{shabat}, a comunidade judaica é
unificada em um espaço de pertencimento e de identificação, tornando
esta territorialidade um espaço sagrado, ``Seu Mundo'', um local
que permite o descanso (e o encontro com a paz) e uma territorialidade
de presença e de devoção divina. Como menciona a escritora bell hooks:

\begin{quote}
Os espaços podem ser reais e imaginários. Os espaços podem contar
histórias e desdobrar relatos históricos. Os espaços podem ser
interrompidos, apropriados e transformados pela prática artística e
literária\footnote{\textsc{hooks}; \textsc{yearning}. \emph{Race, gender and cultural
  politics}, p.\,23. Tradução da autora.}.
\end{quote}

Enquanto um artifício, um objeto e um território temporário, e assim,
diaspórico, o \emph{eruv} não acontece, nem mesmo pertence, a um
determinado espaço, já que ``um espaço se torna um lugar quando se
coloca uma série de objetos nele''\footnote{\textsc{levy}, \emph{op.\,cit.}, p.\,59. Tradução da autora.}. O artifício do \emph{eruv} parte de seu
agenciamento, e a ``agência dos objetos'' permite que estes artifícios
se tornem suporte para algo invisível.

O \emph{eruv} não é um não lugar, ou melhor, o artifício do \emph{eruv}
não se estabelece enquanto uma utopia. O \emph{eruv} é um território
real e físico --- um espaço-tempo apropriado pelas comunidades judaicas
ortodoxas ---, em que outros tempos, outras frequências e outras
relações de sociabilidade se fixam, se concretizam, manifestam-se para
que, então, desmanchem-se\footnote{Desmanchar é o ato de alterar ou
  desfazer o aspecto, a forma ou a arrumação de um determinado objeto,
  bem como desarranjar, desalinhar e desarrumar suas estruturas.
  ``Desmanchar'' é o verbo que permite que partículas moleculares se
  desintegrem. O objeto/\,corpo simplesmente some e seus resquícios ---
  que passam a ser microscópicos --- voam pelo ar. O \emph{eruv}, em seu
  aspecto visível-invisível, se desintegra com facilidade. Por ser um
  objeto móvel e diaspórico permite que seus limites se expandam (ou se
  retraiam), construindo e consolidando localidades --- até que estes
  locais se desvaneçam.}. O \emph{eruv} é uma heterotopia
territorial\footnote{Heterotopia é um conceito da geografia humana que
  descreve lugares que funcionam em condições não hegemônicas, com
  múltiplas camadas de significação. As heterotopias são espaços de
  alteridade, simultaneamente físicos e mentais.} que permite pensar os
espaços (e os tempos) como multiformes --- um interminável e simultâneo
cruzamento de localidades, de tempos e de situações.

Assim, o \emph{eruv} pode ser analisado enquanto um signo da diáspora,
já que se estabelece enquanto uma construção real e fictícia situada em
um determinado espaço urbano em que a comunidade decide fixar-se. O
\emph{eruv} são fios e postes móveis, que se expandem ou se encolhem a
depender das necessidades, dos laços e dos traços de sociabilidade. O
\emph{eruv} é uma fronteira fluida, uma territorialidade elástica de
apropriação e de pertencimento --- um limite e um limiar em si.

\begin{center}\adforn{49}\end{center}

\noindent{}O \emph{eruv} é um fio de diâmetro milimétrico, transparente. Uma
demarcação e uma delimitação espacial judaica real e imaginária,
concreta e abstrata, visível e invisível, física e ficcional, literária
e literal, legal e comunitária. O \emph{eruv} é uma geografia de
identidade temporária e apropriada. São como ilhas em meio ao oceano
que, antes mesmo de se mostrarem, desaparecem --- para que então possam
nascer em outras localidades e em outros tempos.

Em momentos de diáspora estamos todos nós em busca de um refúgio. Como
menciona o filósofo Denètém Touam Bona\footnote{\textsc{bona}, \emph{op.\,cit.}, p.\,52.}:
``destinados a nos tornar refugiados em nossas próprias terras,
todos estamos em busca de um solo, um ar, uma água que não tenham sido
corrompidos, estamos todos em busca de um fora''\footnote{\emph{Ibidem}.
  Denètém Touam Bona, filho de pai centro-africano e de mãe francesa,
  procura construir pontes e sentidos para uma identidade fronteiriça em
  um mundo que, ainda hoje, torce a ``linha da cor''. Em seus
  projetos, obras e escritos, o autor busca fazer do aquilombamento (a
  arte da fuga e evasão dos escravizados) um objeto filosófico por si
  só, usado para pensar o mundo contemporâneo, as fronteiras, os
  corpos/\,comunidades desviantes, bem como as constantes necessidades de
  fuga --- de territorialidades físicas ou imaginárias.}. Em uma
perspectiva diaspórica, o \emph{eruv} é uma espécie de fora, de outras
temporalidades e territorialidades de agregação e de apropriação que
surgem para além de um espaço (e de um tempo) físico homogêneo e
hegemônico que nos é constantemente apresentado. O \emph{eruv} é uma
fronteira e, assim como todo limite, ``opera a inscrição espacial
de uma coletividade humana em um dado território''\footnote{\emph{Ibidem}, p.
  42.}, simplesmente por instaurar e delimitar este espaço. Como nos
lembra Glória Anzalduá, as fronteiras são esses espaços-territórios de
encontro com o outro --- sejam eles geográficos ou linguísticos.

No livro-ensaio \emph{TAZ}, o filósofo e escritor Hakim Bey sustenta que
as chamadas ``zonas autônomas temporárias'' são táticas
sociopolíticas para criar espaços temporários que escapam às estruturas
formais de controle. Estas outras frequências, que se firmam enquanto
possibilidades de reinventar e reivindicar territórios, são
``operações de guerrilha'' que têm por intuito a libertação de
uma determinada área de terra, de tempo (e de imaginação) por um
deliberado momento.

O \emph{eruv} é um signo de pertencimento, um território de agregação
simbólico e diaspórico, coletivo e comunitário --- um espaço
essencialmente móvel e temporário. Nesta perspectiva, o \emph{eruv} pode
ser compreendido enquanto uma ``zona autônoma temporária'' na
medida em que se fixa, para logo dissolver-se e se refazer em outros
territórios e em outros tempos.

Os espaços, as territorialidades, os locais de pertencimento e de
apropriação --- sejam eles físicos, comunitários ou imaginários --- se
mostram enquanto um conglomerado de tempos e de situações. Nestas
localidades o tempo em si não é linear, ou cíclico/\,elipsoidal, é antes
um ``rizoma temporal''\footnote{\textsc{rago}, \emph{op.\,cit.}, p.\,13.}, ou seja,
um espaço heterotópico em que há o acréscimo de múltiplos e simultâneos
tempos em uma só localidade\footnote{O tempo não é elipsoidal, muito
  menos linear. Trata-se de pensar o tempo enquanto uma multiplicidade:
``coexistência de múltiplas temporalidades em um mesmo momento,
  objeto, situação, relação''. Trata-se, aqui, de pensar o tempo como
  um rizoma temporal. \emph{Ibidem}, p.\,11.}. As heterotopias são alternativas para a
contestação dos espaços existentes, uma possibilidade palpável de
reinventarmos e de darmos novos sentidos às espacialidades físicas,
geográficas, políticas, afetivas e subjetivas, já que ``no seio de
uma heterotopia existe uma ruptura do tempo real''\footnote{\emph{Ibidem}, p.\,14.}.

Porém, ``ao invés de nos deslocarmos para um lugar irreal, para um
espaço e tempo imaginários, seria mais interessante nos voltarmos para o
que nos circunda, estranharmos o que é familiar, percebê-lo
diferentemente e reinventá-lo''\footnote{\textsc{foucault} \emph{apud} ibidem, p.\,14.}.
Os espaços heterotópicos inquietam pois dizem sobre o ``aqui'' e
o ``agora'' --- uma possibilidade de transformar o mundo interior
e exterior, individual e coletivamente. São nestes espaços-tempos
heterotópicos, múltiplos e simultâneos que se dão outros sentidos para
as agregações comunitárias e coletivas.

\chapter{Lugares}

\noindent{}Conta Martin Buber: ``Perguntaram ao Rabi Pinkhas: Por que é que,
como afirma a tradição, o Messias deverá nascer no dia do aniversário da
destruição do Templo? O grão semeado na terra --- replicou ele --- deve
apodrecer, para que brote a nova espiga. O poder não pode ressuscitar,
sem antes dissolver-se no grande segredo. Despir uma forma, assumir uma
forma, isso se faz no instante do puro nada. Na concha do esquecimento,
cresce o poder da memória. No dia da destruição, o poder jaz no fundo e
aí cresce. É por isso que nesse dia nos sentamos no chão. É por isso que
nesse dia visitamos os túmulos. É por isso que nesse dia há de nascer o
Messias''\footnote{\textsc{buber}, \emph{op.\,cit.}, p.166.}.

Os casamentos judaicos podem ser realizados em quaisquer lugares,
contanto que sejam feitos debaixo de uma \emph{chupá}: um pedaço de pano
retangular ou quadrado tensionado, formando uma espécie de cabana, e
sustentado por quatro pilares, cada qual em uma extremidade. A
\emph{chupá} é um espaço sagrado em que o casal se encontra para selar a
união entre si e com Deus. Finaliza-se este ritual com a quebra de um
copo de vidro sob o chão pisado pelo par. Simultaneamente, ou seja, no
exato instante desta ruptura, o pedaço de pano --- que representa o
lugar santo da \emph{chupá} --- é recolhido. Já não há mais espaço
sagrado, o copo foi quebrado e a \emph{Shechiná} --- a habitação e a
presença de Deus que estava dentro do território temporário e apropriado
da \emph{chupá} --- foi, em forma de vento e em forma de sopro,
espalhada pelo universo.

O espraiamento e a dispersão\footnote{A palavra ``dispersão'' provém da
  raiz grega ``diasperein'': dispersar, espalhar, disseminar. O prefixo
  ``dia'' (através de) unido à raiz ``sperein'' (esporos ou sementes)
  revela, ao final, que ``diasperein'' significa através de sementes ou
  de esporos. Dessa forma, toda dispersão é uma semeadura, ou seja, tudo
  nasce para que seja espalhado/\,dispersado e levado pelo velho vento do
  sopro nômade da diáspora para que semeie e que brote em outras
  localidades.} da \emph{Shechiná} pelo mundo se apresentam enquanto um
momento de uma lamentação tão profunda que se transforma em festividade.
Assim como é, também, a música iídiche. O estilo
\emph{klezmer}\footnote{\emph{Klezmer} é uma palavra de origem
  iídiche que significa instrumentos musicais, e é também um
  gênero musical não litúrgico desenvolvido em meados do século \textsc{xv} pela
  cultura \emph{ashkenazi} (judeus provenientes da Europa Central e do
  Leste Europeu), em que a dança é um elemento determinante e as
  melodias imitam os risos, choros e suspiros dos \emph{hassidim}. \emph{Hassidim} (plural de \emph{hassid}) são os membros do
  \emph{hassidismo} judaico, um movimento que dá ênfase à alegria
  religiosa e é organizado por dinastias, tendo sempre um líder
  espiritual em cada uma das gerações.} são melodias, cantos e cantigas
de tanta melancolia e lamentação que se transformam em espaços de
celebração de vida.

A quebra do copo de vidro é uma recordação de que judeus e judias estão
incompletos sem a presença do Templo Sagrado de Jerusalém. Porém,
deve-se ter em mente que todos os eventos evocadores do Templo se firmam
enquanto momentos de celebração, assim como é o casamento.

Anteriormente à quebra do copo lê-se o Salmo 137, que diz: 

\begin{verse}
Se eu me esquecer de Ti, ó Jerusalém\\
Que minha mão direita resseque\\
Que minha língua cole no céu da boca\\
Se eu não lembrar de Ti\\
Se eu não fizer de Ti, ó Jerusalém\\
a maior das minhas alegrias.
\end{verse}

Este versículo é uma afirmação de devoção e de vida das comunidades judaicas, onde
constantemente se rememora a incompletude com a ausência do Templo
Sagrado, mesmo nos momentos de maior felicidade. Assim, como diz Rabi
Pinkhas: ``{[}\ldots{}{]} No dia da destruição o poder jaz no fundo e
aí cresce''\footnote{\textsc{buber}, \emph{op.\,cit.}, p.166.}. É nestes momentos de
ruptura, de quebra, de lamentação e de melancolia que os povos em
diáspora, antes de qualquer coisa, celebram a vida.

Comentam os sábios que a santidade espiritual atingida no momento da
\emph{chupá} é tamanha que se encontram neste local sagrado --- e
temporário --- as dez últimas gerações de pessoas falecidas que fazem
parte destes fios, destes nós e destas linhas genealógicas. As almas
descem (e descendem\footnote{Se partirmos do pressuposto de que tudo vem
  do pó, para retornar ao pó, então as almas retornam, descem, e ao
  mesmo tempo descendem, da esfera divina. A vida e a morte estão no
  mesmo lugar.}) da esfera divina e agregam-se no local da \emph{chupá}
abençoando e santificando a união do casal.

Por fim, a \emph{chupá} cria um lugar santo e temporário, transitório e
apropriado em que se encontram, ao mesmo tempo, Deus e toda a sua
criação. Deus é toda a criação (\emph{HaMakom}\footnote{\emph{HaMakom}
  significa, em português, ``o Lugar'' e é um dos nomes judaicos de
  Deus.}), é o lugar por completo que engloba os vivos e os mortos. Deus
envolve não somente os lugares em que estamos --- em que vivemos ---,
mas toda a existência em si. E \emph{HaMakom} envolve tudo, assim como o
útero envolve o feto.

Durante a santidade do momento do \emph{shabat} clama-se: ``Vem,
meu Amado, saudar a Noiva, vem, vamos dar as boas-vindas à Presença do
shabat''. A canção de \emph{Lechá Dodi}\footnote{O \emph{Lechá Dodi} é
  um poema cantado e entoado no momento de início do \emph{shabat}. Esta
  cantiga é de autoria do Rabi Shlomo Halevi Alkabets, composta entre os
  anos de 1529 e 1540. Porém foi Isaac Luria (1534--1572), conhecido como
  Ari, que desempenhou um papel fundamental para a popularização e
  espiritualização do hino.} é o ponto máximo da elevação espiritual do
\emph{cabalat shabat} (literalmente ``o recebimento do \emph{shabat}''), realizado às
sextas-feiras à noite. Nesse momento, e por meio desta bênção, Deus é
convidado a unir-se com os humanos e habitar a esfera do \emph{shabat}.
A noiva é o próprio \emph{shabat}, pois, como nos diz o
\emph{midrash}\footnote{O \emph{midrash} é uma metodologia judaica de
  ler as escrituras sagradas. Seu fundamento está em uma visão cíclica
  da história, em que o acontecimento central, ou melhor, a presença de
  Deus no texto, é continuamente revisado, oferecendo sempre novos
  significados às várias situações expostas, bem como às interpretações.}:
``Deus criou o tempo, bem como o dividiu em sete dias, e prometeu
ao \emph{shabat} (o sétimo dia da criação) que \emph{Israel será seu novo
par}\footnote{\emph{Midrash: Bereshit Rabá} 11.} \emph{para toda a
eternidade''.} Assim, todas as semanas o povo judeu saúda e celebra a
aproximação do \emph{shabat}, como se fosse um noivo que aguarda a
chegada de sua noiva.

Segundo a interpretação de Anat Yosef,\footnote{Anat Yosef foi um
  intérprete midráshico, isto é, comentarista das escrituras bíblicas.}
na verdade, a ``noiva'' mencionada pode ser entendida enquanto uma
alusão à própria \emph{Shechiná} --- a habitação e a presença de Deus.
Ou seja, durante o \emph{shabat} oramos pela aproximação da presença de
Deus no mundo humano. Por fim, segundo a filosofia judaica, se a
\emph{Shechiná} é a ``noiva'' --- que se aproxima do \emph{shabat} --- e
o povo judeu é o ``noivo'', então estamos a todo instante na esfera de
Deus, clamando pela presença de Sua habitação. Estamos em Deus, e
instamos para que a Sua existência esteja em nós também.

O Rabi Kobrin ensinava: ``Deus fala ao homem como falou a Moisés:
\emph{Tira os teus sapatos dos teus pés, descalça o habitual que encerra teus
pés e saberás que o lugar em que pisas é solo sagrado. Pois não há
situação da vida humana em que não se possa encontrar, em toda a parte e
a qualquer tempo, a santidade do Senhor}''\footnote{\textsc{buber}, \emph{op.\,cit.},
  p.\,490.}.

Por fim, se o universo --- e todas as coisas\footnote{Neste ensaio
  utiliza-se o sentido heideggeriano de ``coisas'', ou seja, as coisas
  como encontros de fios e de vida.} que ele contém, sejam elas vivas ou
mortas --- está dentro da esfera divina, então a esfera de Deus engloba
tanto o \emph{Olam Hazé} (``este mundo''), quanto o \emph{Olam Habá}
(``o mundo vindouro''), lembrando que a palavra hebraica \emph{Olam} é
usada tanto para designar ``tempo'', quanto para designar ``espaço''. E
Deus é o tempo e o espaço em que tudo acontece --- o lugar, a
completude, o eterno. Como nos lembra o pesquisador e estudioso das
escrituras bíblicas Jeff Benner:

\begin{quote}
Palavras hebraicas usadas para o espaço também são usadas para o tempo.
{[}\ldots{}{]} A palavra hebraica Olam significa literalmente ``além do
horizonte''. Ao olhar para longe, é difícil perceber detalhes e o que
está além desse horizonte não pode ser visto. Este conceito é o Olam. A
palavra Olam também é usada para o tempo: para o passado distante ou
para o futuro distante, como um tempo difícil de conhecer ou perceber.
Esta palavra é frequentemente traduzida como ``eternidade'', significa
um período de contínuo que nunca termina\footnote{BENNER,
  \emph{Eternity}.}.
\end{quote}

Assim, Deus não é infinito, mas eterno. Não há nem começo, nem fim em
Deus --- tudo é uma questão de meios (como são também as escrituras e as
interpretações do Talmud, ou as pedras dispostas sobre os
túmulos judaicos). O Deus judaico é a completude que habita os eventos e
os movimentos, e o judaísmo, então, é uma concepção do tempo.

Voltemos ao Rabi Pinkhas: ``{[}\ldots{}{]} O grão semeado na terra
{[}\ldots{}{]} deve apodrecer, para que brote a nova espiga''\footnote{\textsc{buber},
  \emph{op.\,cit.}, p.166.}. Estes fragmentos de grãos, que são --- ao mesmo
tempo --- fragmentos de vida e de morte, são também sementes de
celebração. Estas sementes hão de morrer para renascer uma nova espiga.
E assim, Didi-Huberman nos lembra que estas sementes de renascimento de
vida devem ser guardadas ``não como tesouros imutáveis, mas como
sementes para o presente, para o futuro''\footnote{\textsc{didi-huberman}, op.
  cit., p.\,129.}.

Os poderes germinativos das sementes --- como aquelas encontradas nos
papéis do gueto de Varsóvia --- não foram armazenados para dar soluções
àqueles que as coletaram (essas pessoas já são cinzas\footnote{Estes
  indivíduos não sobreviveram, mas suas histórias sim.}). Ao contrário,
as sementes são respostas (não conclusivas) de vida para outrem, para o
presente e para o futuro --- para nós, e para a nossa geração. Por fim,
como nos lembra Rabi Pinkhas, se tudo nasce e tudo morre, para então
renascer, ou melhor, se ``{[}\ldots{}{]} No dia da destruição, o poder
jaz no fundo e aí cresce''\footnote{\textsc{buber}, \emph{op.\,cit.}, p.166.}
então, ``{[}\ldots{}{]} É por isso que nesse dia há de nascer o
Messias''\footnote{\emph{Ibidem}.}.

Entendendo que o ciclo nunca se completa sempre da mesma maneira, que
Deus habita o movimento\footnote{Deus habita os textos, os livros e o
  próprio trânsito da diáspora.} e ``os judeus, como tal, rejeitam
a estaticidade das coisas e das ideias, e acreditam na transformação e
na reedição''\footnote{\textsc{zevi}, \emph{Arquitetura e Judaísmo: Mendelsohn},
  p.\,10.}, nos permitimos aqui revisitar as nossas histórias --- de
judeus ou não judeus. Que possamos nos apropriar delas, da morte para o
renascimento, e da vida enquanto um princípio contínuo.

\begin{center}\adforn{49}\end{center}

\noindent{}``Quando o Rabi Abraão Iaakov, filho do Rabi Israel de Rijin,
ainda era criança, passeava ele no quintal, para cima e para baixo, ao
anoitecer de sexta-feira, quando os hassidim já haviam ido rezar. Um
hassid dirigiu-se a ele e falou: \emph{Por que não entras? Já é \emph{shabat}}.
\emph{Ainda não é} --- respondeu Abraão. \emph{Como sabes isso?} --- perguntou o
hassid. \emph{No \emph{shabat}} --- replicou o menino --- \emph{surge sempre um novo céu
e ainda não o estou vendo}\footnote{\textsc{buber}, \emph{op.\,cit.}, p.\,383.}.

O calendário (e o seu tempo) é essencialmente lunar. Dessa forma, os
dias da semana, o \emph{shabat} e todos os feriados, as festas e as
festividades judaicas iniciam-se sempre no período de crepúsculo. Ao
final, o anoitecer permite um ``novo céu''\footnote{\emph{Ibidem}.}, um
novo recomeço. O dia judaico inicia ao nascer das três primeiras
estrelas no céu e termina no momento do pôr do sol.

As interpretações rabínicas sobre os textos sagrados não possuem
definições claras e concretas a respeito do momento do ``anoitecer''.
Porém, segundo o Talmud e o \emph{Midrash}, é nesse instante que
Deus, antes da entrega do \emph{shabat}, ao final do sexto dia da
criação divina, cria dez elementos, sendo eles: o arco-íris, o
\emph{maná}\footnote{\emph{Maná} é um alimento produzido milagrosamente,
  uma dádiva divina aos hebreus, durante toda a peregrinação no deserto
  do \emph{Sinai}, rumo à terra de \emph{Canaã}.}, a boca da terra, a
boca do jumento, as Tábuas da Lei, as letras, a escrita, o túmulo de
Moisés, o poço e o \emph{shamir}\footnote{Conta a tradição que o
  \emph{shamir} é um ser vivo criado por Deus e dotado da capacidade de
  cortar qualquer tipo de material, inclusive a mais dura rocha. O
  \emph{shamir} foi utilizado para a construção do Templo de Jerusalém,
  pois Deus não permitiria a utilização de arsenais de guerra para a
  construção deste lugar sagrado.}.

Entre o momento do pôr do sol e a aparição das três primeiras estrelas
na esfera celeste, os judeus se encontram em um limiar. Para a filosofia
judaica, nesse intervalo (e assim, nesse espaço-situação), não é nem dia
(em hebraico: \emph{iom}), nem noite (em hebraico: \emph{laila}), mas
sim um espaço-tempo entre o dia e a noite, chamado de \emph{erev} (em
português: ``tarde'').

\emph{Erev} é uma palavra hebraica cognata e da mesma raiz que as
palavras \emph{eruv} e \emph{arab}. \emph{Arab} significa ``mistura'',
\emph{eruv} é uma combinação de pátios e \emph{erev} é uma mescla entre
tempos, entre o dia e a noite. Esse limiar é conhecido como \emph{bein
hashemashot}, em português, ``entre sóis''.

Com uma duração de aproximadamente dezoito minutos (a depender da
localização geográfica), essa situação-período permite, por exemplo, o
acendimento das velas de \emph{shabat}\footnote{Por conta do judaísmo
  possuir um calendário lunar, todas as festas, festividades e
  celebrações judaicas iniciam-se no período do anoitecer, por meio do
  acendimento de fogo.}, em um espaço-tempo que não é ainda o
\emph{shabat}, e muito menos o sexto dia da criação divina, mas a
passagem entre estes dois períodos. Assim, \emph{bein hashemashot} é um
período, um espaço e um tempo limite e um limiar em si. Um momento entre
o dia e a noite, nem aqui, nem lá --- trânsito.

Como comenta Rabi Pinkhas: ``Despir uma forma, assumir uma forma,
isso se faz no instante do puro nada''\footnote{\textsc{buber}, \emph{op.\,cit.}, p.166.}.
A passagem entre o dia e a noite (a mudança e a transformação de uma
situação) se firma enquanto uma transição de forma, de matéria e de
espírito --- sair de um espaço e de um momento para entrar em outro.
Esta transição ou rito de passagem necessita de um tempo de pausa, de
reflexão, renovação e assim, de renascimento. Este movimento se faz no
``instante do puro nada''\footnote{\emph{Ibidem}.} --- no limiar e na
ambiguidade de conotações e de definições. O ``instante do puro
nada''\footnote{\emph{Ibidem}.} não é o vazio. O nada é a possibilidade.

\emph{Twilight People}, poema litúrgico do rabino norte-americano Reuven
Zellmann, se firma enquanto um lembrete de que este momento do
crepúsculo é um espaço plural e coletivo para todos. Tal espaço-tempo
sagrado de intermédio, instante de transformação e de renascimento,
reaparece em diversas religiões e crenças ao redor do mundo.

No poema do rabino Zellmann abençoa-se essa qualidade crepuscular dos
seres humanos. Abençoa-se a criação de um espaço-tempo de intermédio em
que as relações são interrompidas e o binarismo já não existe.
Abençoa-se o momento em que as crenças e as certezas absolutas são
suspensas e o julgamento é suavizado e esvaziado. Abençoa-se uma
situação de possibilidade, um momento do nada. E, finalmente, abençoa-se
Deus --- o criador do Universo, do dia e da noite, dos ventos, das
nuvens e do crepúsculo --- Aquele que transcende todas as categorias e
definições.

Entende-se aqui que a diáspora, ou os lugares judaicos, são espaços de
intermédio, que se firmam enquanto caminhos entrecruzados de memórias e
de destinos, de laços e de nós. A diáspora não pode ser determinada, mas
sempre pode ser vivida.

\begin{center}\adforn{49}\end{center}

\noindent{}Em \emph{Borderland/\,La Frontera}, a poeta e ensaísta Gloria Anzaldúa diz
que os grupos fronteiriços vivem no entre, em \emph{nepantla}:
``nepantla é estar em meio a uma transformação {[}\ldots{}{]} nepantla
é um jeito de ler o mundo''\footnote{\textsc{anzaldúa}, \emph{op.\,cit.}, p.\,237. Tradução da autora.}. Viver no entre --- no limiar --- é uma condição das
comunidades e da identidade judaica. Estar em fronteira é estar
diretamente ligado à alteridade, em contato e conexão com os demais
grupos sociais que se estabelecem nestes territórios.

Em \emph{Limiar, aura e rememoração}, a pesquisadora Jeanne Marie
Gagnebin argumenta que estes limiares são zonas de intermédio e rituais
de passagem que permitem suspender o espaço, e assim fugir de definições
conclusivas e autocircundantes, tais como: público/\,privado, vida/\,morte,
sagrado/\,profano. A experiência do limiar pode ser lida enquanto uma
soleira ou o umbral de uma porta (da mesma maneira que é o \emph{eruv}),
uma zona de transição e de mudança.

Atravessando fronteiras, os povos em diáspora contestam e questionam as
conotações. E, dessa forma, o viver no limiar permite, justamente, a
abertura de poros para o movimento de não apenas ocupar, mas sim de
habitar o entre. Este espaço entre, finalmente, possibilita
``reconquistar para o pensamento os territórios do indeterminado e
do intermediário, da suspensão e da hesitação''\footnote{GAGNEBIN,
  \emph{Limiar, aura e rememoração}, p.\,40.}.

Ao habitar a diáspora, a identidade judaica e as agregações
comunitárias, por meio de suas tradições, cultos e ritos, bem como de
seus idiomas e de suas línguas, se apropriam dos espaços de trânsito ---
estes movimentos de dispersão cíclicos e constantes --- e os tornam uma
possibilidade. Em outras palavras, as comunidades judaicas (e seus
indivíduos) transformam essas condições de existência --- estar entre
espaços --- em um local (e um tempo) de pertencimento.

Os movimentos apropriativos da condição diaspórica assemelham-se a
ecótonos. Ecótonos são zonas de transição ambiental, regiões resultantes
do contato entre dois ou mais biomas fronteiriços, em que se encontram
as mais diversas comunidades ecológicas. São territórios ricos em
espécies, pois compartilham espaços e permitem que sejam influenciados
por outros ambientes. É exatamente desta maneira que devem ser lidas as
identidades fronteiriças. A possibilidade de ``não estar nem aqui, nem
ali'' permite o contato, a troca e a multiplicidade. A diáspora é uma
condição em si.

Por fim, o antropólogo e pesquisador Tim Ingold, em \emph{Estar Vivo},
nos lembra que a vida nunca é, exclusivamente, vivida neste ou naquele
lugar, ``mas sempre no caminho de um lugar ao outro''\footnote{INGOLD,
  \emph{Estar vivo: ensaios sobre movimento, conhecimento e descrição},
  p.\,218.}. A vida não é vivida dentro de um espaço, mas através dele
--- na fronteira.

Essa possibilidade de viver na fronteira, e na diáspora, faz com que os
grupos sociais habitem a terra como peregrinos, em transição, em
deslocamentos e em constantes mudanças, pois ``a experiência humana
não é situada, mas situante {[}\ldots{}{]} Ela desdobra-se não em lugares,
mas ao longo de caminhos''\footnote{\emph{Ibidem}, p.\,219.}. Assim, os
peregrinos caminham, física e mentalmente, e tecem linhas, traços, laços
e nós coletivos.

Em \emph{Estar vivo}, Tim Ingold menciona o livro \emph{Playing Dead} do
escritor canadense Rudy Wiede, que fala do povo \emph{Inuit}, habitante
da região ártica do Canadá, Alasca e Groenlândia. Ingold argumenta que,
para os \emph{Inuit}, ``assim que uma pessoa se move, ela se torna
uma linha''\footnote{\emph{Ibidem}, p.\,221.}. As linhas se transformam em
caminhos pelos quais a vida é vivida (ou escrita, como no caso dos
papéis do gueto), pois o peregrino está em contínuo movimento,
justamente por ele próprio ser o movimento. O peregrino ``é
exemplificado no mundo como uma linha de viagem {[}\ldots{}{]} É uma linha
que avança a partir da extremidade conforme ele segue
adiante''\footnote{\emph{Ibidem}.}. É por meio destes movimentos físicos, ou
destes movimentos mentais, que as linhas consolidam-se enquanto uma
possibilidade de habitar o deslocamento.

% Comentário da editora: A tradução da frase do peregrino está estranha. Ela está correta?

Ao compreender as identidades diaspóricas enquanto peregrinos que
habitam a terra, entende-se aqui que estes indivíduos ``não têm
destino final, pois onde quer que estejam, e enquanto a sua vida
perdure, há algum outro lugar aonde podem ir''\footnote{\emph{Ibidem}.}.
% Comentário da editora: é "terra" mesmo, né?

Existem duas palavras hebraicas para designar o exílio: \emph{galut} e
\emph{tefutzá}. \emph{Galut} (em português: exílio) é visto como algo
maldito, sofrido e a ser superado. \emph{Tefutzá} (em português:
dispersão), em contrapartida, é visto como um lugar em si. O exilado se
perdeu, o disperso está perdido. Ao final, o movimento de ``estar
perdido'' é uma posição política, um ato de possibilidade, de viver o
entre, o nada --- habitar a fronteira.

\begin{center}\adforn{49}\end{center}

\noindent{}Em ``Contra o espaço: lugar, movimento, conhecimento'', Tim Ingold diz
que ``espaço'' é um dos termos mais abstratos para descrever o mundo em
que vivemos, na medida em que pensar que a vida humana acontece em um
espaço --- ou seja, que ocupa uma determinada espacialidade --- está
diretamente vinculado ao que o antropólogo chama de ``lógica da
inversão''. Segundo o pesquisador, a ``lógica da inversão''
transforma as vias ao longo das quais a vida é vivida em limites nos
quais a vida é encerrada.

A diáspora é o mundo habitado, e não o mundo ocupado. A diáspora tece
lógicas a partir da construção de nós e de fios do devir, ou seja, a
partir das peregrinações e dos deslocamentos, de indivíduos, de
culturas, de tradições, de identidades, bem como de comunidades. Os
``humanos habitam a terra como peregrinos''\footnote{\emph{Ibidem}, p.
  219.} --- semeiam e cultivam nós, até que estes laços brotem para logo
se desmancharem, e serem refeitos em outras localidades. Por assim
dizer, ``lugares, então, são como nós, e os fios a partir dos
quais são atados linhas de peregrinação''\footnote{\emph{Ibidem}, p.\,220.}.

Se os movimentos diaspóricos judaicos se apropriam dos territórios do
indeterminado por meio de sua cosmogonização (os transformam em casas,
em um lugar habitável, e não espaço ocupado), e se o Deus judaico habita
o movimento, então há uma necessidade judaica de contínua e
recorrentemente se desvincular da ``abstração do
espaço''\footnote{\emph{Ibidem}, p.\,216.} para, então, habitar o lugar.

É a partir da inserção de objetos, de símbolos e signos de pertencimento
e de identificação que as comunidades judaicas criam o sentimento de
``lugar''. Há uma constante fuga do ``sentimento abstrato do
espaço''\footnote{\emph{Ibidem}.}, já que, quanto mais alta a dimensão do
espaço, mais distante o embasamento no lugar --- no ``Seu Mundo''.

Ingold argumenta que cada lugar contém, em si, lugares menores.
``Os lugares sempre se abrem para revelar outros lugares dentro de
si''\footnote{\emph{Ibidem}.}, como bonecas russas ou espaços heterotópicos.
Os lugares existem no espaço. E os lugares judaicos habitam o ``tempo'',
pois fazem do tempo o seu lugar.

O rabino e filósofo Abraham Heschel diz que ``os sábados são
nossas grandes catedrais''. O momento judaico do \emph{shabat} é um
lugar, justamente pela experiência do \emph{shabat} ser totalmente
desvinculada de um espaço. Quanto mais próximos os lugares, mais estão
atados as linhas e os laços entre a comunidade. O \emph{shabat} é
``uma volta ao lar''\footnote{\textsc{buber}, \emph{op.\,cit.}, p.\,196.}
justamente por cristalizar uma concepção de lugar.

Em diáspora, a experiência judaica vive em múltiplos e multifacetados
lugares ao mesmo tempo: ``a eternidade é conquistada por aqueles
que trocam o espaço pelo tempo''\footnote{\textsc{zevi}, \emph{op.\,cit.}, p.\,11.}.

\begin{center}\adforn{49}\end{center}

\noindent{}Os pisos das sinagogas sefaraditas\footnote{Sefaraditas
  são as comunidades judaicas provenientes de Portugal e da Espanha e
  que possuem a língua diaspórica ladino como idioma mãe.}
encontradas no Caribe (em especial no Suriname, em Curaçao e na Jamaica)
são completamente cobertos com areia proveniente de
\emph{Sefarad}\footnote{\emph{Sefarad} é o termo hebraico correspondente
  à região atual da Espanha.}. Estas sinagogas foram construídas em
meados do século \textsc{xvi}, durante o período da Inquisição Espanhola que se
alastrou por Portugal, bem como por suas colônias, e desenraizou os
judeus de seus países de origem.

Em meio à perseguição, a areia despejada sobre o piso permitia um
abafamento dos ruídos dos passos dos membros da congregação, e assim
fazia com que a vida judaica pudesse ser vivida em silêncio e em
segredo.

Para além de um movimento de resistência e de contestação, o ato de
esparramar areia sobre o piso sinagogal pode ser lido também enquanto um
lembrete de devoção à divindade e de memória ao momento de vagueio por
quarenta anos no deserto do \emph{Sinai}, antes da entrada na terra de
\emph{Canaã}. Mais que isso, a areia de \emph{Sefarad} implica uma
recordação das relações de intimidade que as comunidades
sefaraditas possuem com a sua terra natal (a Espanha) e com a
diáspora. E, afinal, o que se leva de um lugar ao outro?

O pesquisador Adolfo Roitman afirma que ``todo lugar pode ser
prima facie uma sinagoga, implicando que o elo com Deus não reside em um
lugar fixo e determinado (como é o caso de um templo), mas que se fazia
presente onde Israel se reunia e elevava suas preces''\footnote{\textsc{roitman},
  \emph{Del Tabernáculo al Templo}, p.\,59.}. Assim sendo, a tradição
judaica entende que são os indivíduos que incorporam em si e entre si os
lugares de pertencimento ao judaísmo. Ou seja, o ``ser diaspórico'' é um
movimento compartilhado, um saber ancestral judaico e coletivo, uma
memória que acompanha as comunidades.

Estes constantes movimentos de dispersão e de transitoriedade permitem
que as comunidades judaicas e cada um dos que a compõem sejam peregrinos
e caminhem pelos territórios levando consigo seus lugares sagrados, o
``Seu Mundo'': Deus, os ritos, os símbolos e os objetos hierofânicos.
Todo novo local é demarcado e cosmogonizado.

A escritora Noemi Jaffe diz que ``ser judia é, de alguma forma,
continuar a costurar''\footnote{\textsc{jaffe}, \emph{Rebeca: A costura no
  avesso}.}. Que continuemos a tecer estes caminhos herdados e
compartilhados. Caminhos que, em algum sentido, já foram percorridos, e
serão mais uma vez, por nós. Mais uma vez apropriados, questionados,
relidos. Os caminhos judaicos estão em constante \emph{chidush}. Todos
podem e devem se apropriar, fazer deles um pouco seus.

\begin{center}\adforn{49}\end{center}

\noindent{}A intenção deste ensaio é se apresentar enquanto uma história judaica,
pois ``a vida judaica é medida pelos seus livros e permeada pela
história''\footnote{\textsc{zevi}, \emph{op.\,cit.}, p.\,7.}. Dessa forma, a tendência do
espírito judaico da diáspora é de expressar os fatos e os acontecimentos
do passado e do presente de maneiras aguçadas. As ocorrências são de tal
modo relatadas e vivenciadas que passam a dizer alguma coisa, culminando
em algo realmente dito. ``A anedota judaica afirma-se no relato de
um único incidente que aclara todo um destino. E no incidente único
exprime-se o significado da vida''\footnote{\textsc{buber}, \emph{op.\,cit.}, p.\,15.}.

E assim, na constelação destas páginas, é hora de retornar ao espelho da
página em branco e voltar a escrever palavras esparsas, signos da
diáspora --- apropriar-se da escrita, das letras e, por fim, fazer
nascer sentidos do indeterminado e de identificação.

\emph{Estos son los biervos de mis abuelas i de mis abuelos, ke ayer
cantaron mis papas, i ke hoy soi io quien los canto}\footnote{Do
  ladino: Essas são as palavras de minhas avós e de meus avôs,
  que ontem contaram aos meus pais, e que hoje sou eu quem conto.}. Hoje
sou eu quem conto, quem esboço, risco e delineio as mais diversas
linhas-palavras que querem a todo custo dizer: amanhã --- para assim
haver amanhã.

% Comentário da editora: Eu pararia aqui, e excluiria o Sheheianu.

\noindent{}\emph{Sheheyanu v'kiyimanu Veihiguiyanu lizman hazé}\footnote{``Bendito
  és Tu, {[}\ldots{}{]}, que nos concedeu a vida, nos sustentou e nos
  permitiu chegar a este momento''.}.
